% Twelve pairwise adjacent faces: a chirotope reduction to Boolean unsatisfiability.
% Build:  pdflatex faces12-en && pdflatex faces12-en
\pdfoutput=1 % arXiv: build with pdfLaTeX
\documentclass[11pt,reqno]{amsart}

\usepackage[T1]{fontenc}
\usepackage[utf8]{inputenc}
\usepackage[english]{babel}
\usepackage{amsmath,amssymb,amsthm}
\usepackage{booktabs}
\usepackage{enumitem}
\usepackage{microtype}
\usepackage[a4paper,margin=3.1cm]{geometry}
\usepackage[colorlinks=true,linkcolor=black,citecolor=black,urlcolor=blue]{hyperref}
\usepackage{url}
\hypersetup{
  pdftitle={Twelve pairwise adjacent faces: a chirotope reduction to Boolean unsatisfiability},
  pdfauthor={Andrei Pirozhkov},
  pdfsubject={A conditional reduction for neighbourly polyhedral surfaces},
  pdfkeywords={polyhedral surfaces, oriented matroids, SAT, LRAT, reproducibility}
}
\emergencystretch=1em

\theoremstyle{plain}
\newtheorem{theorem}{Theorem}[section]
\newtheorem{lemma}[theorem]{Lemma}
\newtheorem{proposition}[theorem]{Proposition}
\newtheorem{corollary}[theorem]{Corollary}
\theoremstyle{definition}
\newtheorem{definition}[theorem]{Definition}
\newtheorem{hypothesis}[theorem]{Hypothesis}
\newtheorem{remark}[theorem]{Remark}
\newtheorem{convention}[theorem]{Convention}

\newcommand{\RR}{\mathbb{R}}
\newcommand{\FF}{\mathbb{F}_2}
\newcommand{\bra}[1]{[#1]}
\newcommand{\sgn}{\operatorname{sgn}}
\newcommand{\Ends}{\operatorname{Ends}}
\newcommand{\Aut}{\operatorname{Aut}}
\newcommand{\link}{\operatorname{lk}}
\newcommand{\relint}{\operatorname{relint}}
\newcommand{\OnB}{\mathrm{OnB}}
\newcommand{\Ord}{\mathrm{Ord}}
\newcommand{\refl}{\mathrm{refl}}
\newcommand{\xor}{\oplus}
\DeclareMathOperator{\dett}{det}

\begin{document}

\title[Twelve pairwise adjacent faces]{Twelve pairwise adjacent faces:\\ a chirotope reduction to Boolean unsatisfiability}

\author[Twelve pairwise adjacent faces]{Andrei Pirozhkov}
\thanks{Email: \href{mailto:pirozhkov.andrey.u@gmail.com}{\texttt{pirozhkov.andrey.u@gmail.com}}.
  ORCID: \href{https://orcid.org/0009-0008-6901-7843}{0009-0008-6901-7843}.}
\date{17 October 2026}

\subjclass[2020]{Primary 52B70; Secondary 52C40, 05C10, 68V15}
\keywords{neighbourly polyhedra, oriented matroids, chirotopes, polyhedral realizability, SAT, verified proof checking}

\begin{abstract}
Can twelve planar faces form a closed surface in three-dimensional space with every
pair sharing an edge? We turn this geometric existence question into explicit finite
proof obligations and provide a workflow for checking a distributed computation.
We require an embedded surface in $\RR^3$. Faces are simple planar disks with distinct supporting planes, and any two
faces intersect in precisely one common edge. Such a surface is trivalent and has
genus $6$. We reduce its existence to the satisfiability of one of $59$ explicit
Boolean formulas, one for each orientable neighbourly map on twelve faces.

The reduction first puts a realization into general position without changing its
combinatorics. The twelve supporting planes and the plane at infinity then define
a uniform rank-$4$ chirotope on $13$ elements. Bracket signs encode necessary
conditions for face simplicity, boundary-crossing parity and local vertex geometry.
A reorientation gauge removes redundant signs. Each resulting formula has
$5\,257$ variables and $52\,177$ clauses.

We give the necessity proofs, recompute the $59$ combinatorial types, and describe
exact geometric controls and a separate reference encoder. We prove that refuting
every leaf of a finite tree of case splits rooted at the unrestricted formula refutes
the formula. We also specify a protocol for producing LRAT refutations of such leaves,
checking them with a formally verified checker, and verifying the computation
independently by rechecking or recomputing its leaves.
Unsatisfiability of all $59$ formulas would exclude a surface in the
stated model. A complete set of refutations has not been obtained, so this
non-existence conclusion remains conditional. Conversely, a satisfying assignment
would not by itself establish geometric realizability.
\end{abstract}

\maketitle
\raggedbottom
\enlargethispage{3pt}

\section{Introduction}\label{sec:intro}

\subsection{The question}
Call a closed surface in $\RR^3$ that is a finite union of planar polygons a \emph{polyhedral
surface}, and call it \emph{neighbourly} when every two of its faces share an edge. The tetrahedron
is neighbourly with four faces. Szilassi's toroidal polyhedron \cite{Szilassi1986} is neighbourly
with seven. These are the only examples reported in \cite{ArsenevaEtAl2024}.
We study the next combinatorially admissible case within the model specified below:

\begin{quote}
\textbf{Question.} Is there a closed embedded polyhedral surface in $\RR^3$ satisfying
\textup{(H1)--(H3)} below, with twelve faces, every two of which share an edge?
\end{quote}

The restriction to twelve is not arbitrary. Pairwise adjacency forces $E=\binom{n}{2}$ edges and, as
we show in Lemma~\ref{lem:trivalent}, every vertex of such a surface has degree three, whence
$V=n(n-1)/3$ and
\[
   g \;=\; \frac{2-\chi}{2} \;=\; \frac{(n-3)(n-4)}{12}.
\]
A closed surface embedded in $\RR^3$ is orientable, so $g$ must be a non-negative integer, which
happens exactly for $n\equiv 0,3,4,7\pmod{12}$. After $4$ and $7$ the next admissible value is
$n=12$, with $V=44$, $E=66$ and genus $6$. That the corresponding \emph{combinatorial} objects exist
in abundance is classical: by the Ringel--Youngs solution of the Heawood problem
\cite{RingelYoungs1968,Ringel1974}, the complete graph $K_n$ embeds in the orientable surface of genus
$\lceil (n-3)(n-4)/12\rceil$, and when the unrounded quotient $(n-3)(n-4)/12$ is an integer --- exactly for
these $n$ --- the embedding is a triangulation. The
question is whether one of those combinatorial objects can be realized by flat faces in space
without self-intersection.

\subsection{What is known}
Dualizing a neighbourly polyhedron with $12$ faces gives a triangulation of the orientable surface of
genus $6$ with $12$ vertices in which every two vertices are joined by an edge; $12$ is the minimum
number of vertices for that surface: a triangulation with $v$ vertices has $3(v+10)$ edges,
at most $\binom v2$, which forces $v\ge12$ (see \cite{JungermanRingel1980} for the general
minimum-vertex theorem). Altshuler, Bokowski and Schuchert
\cite{AltshulerBokowskiSchuchert1996} classified these: there are exactly $59$ combinatorial types,
all orientable of genus $6$. They also proved several properties that a polyhedral realization would
have to have, and left realizability open. Schewe \cite{Schewe2010} settled the realizability of the
dual objects negatively: no triangulation of the genus-$6$ surface with $12$ vertices has a
polyhedral realization in $\RR^3$. His Theorem~1.1 follows from Proposition~1.2: none of the $59$
types admits an acyclic uniform rank-$4$ oriented matroid \emph{admissible for that triangulation}.
Admissibility excludes circuits witnessing an intersection of a triangle with a disjoint edge
(\cite[Section~2]{Schewe2010}); a SAT solver discharges this constrained combinatorial problem.

Schewe's theorem does not answer our Question. Polyhedral realizability is not preserved by
combinatorial duality: his objects are $12$ \emph{points} carrying $44$ triangular faces, ours are
$12$ \emph{planes} carrying $44$ trivalent vertices, and a realization of either does not provide a realization
of the other by combinatorial duality alone. The general question --- which complete graphs occur as face-adjacency graphs of closed
polyhedral surfaces --- is recorded as open in \cite{ArsenevaEtAl2024}, where the tetrahedron and the
Szilassi polyhedron are named as the only known instances.

\subsection{Contribution}
This paper gives an explicit reduction of the Question to a finite Boolean computation.
The geometric necessity arguments are proved below; the topological and oriented-matroid
facts used in them are identified by references. Specifically:

\begin{enumerate}[label=(\arabic*),leftmargin=2.2em]
\item \textbf{Case list.} Every neighbourly polyhedron with twelve faces realizes one of exactly $59$
   neighbourly maps (Theorem~\ref{thm:59}). We recompute the list by an exhaustive search
   over oriented triple systems; its size agrees with \cite{AltshulerBokowskiSchuchert1996}, and
   its classes are, one by one, those of the published list \cite{LutzManifoldPage}.
\item \textbf{General position.} Every realization may be perturbed into one in which no four of the
   thirteen covectors (twelve face planes and the plane at infinity) are linearly dependent
   (Theorem~\ref{thm:genpos}). The signs of the $\binom{13}{4}=715$ brackets are then defined and form a
   uniform rank-$4$ chirotope.
\item \textbf{A sign dictionary.} Every geometric predicate the argument needs --- the order of the
   points marked on a line, betweenness inside an edge, the
   turn at a corner, the parity of boundary crossings along a ray --- is an explicit Boolean
   combination of bracket signs (Section~\ref{sec:dictionary}). The identities are exact, not
   approximations.
\item \textbf{Five layers of necessary conditions.} We define a CNF formula $\Phi_T$ for each map $T$
   as the conjunction of five layers $S,H,A,B,I$ and a gauge, and prove that every general-position
   realization of $T$, after the gauge normalization of Corollary~\ref{cor:gauge}, satisfies all of
   them (Theorem~\ref{thm:necessity}). Layers $S$, $H$ and $A$ are
   exactly the statement that the bracket signs form a uniform rank-$4$ chirotope, encoded through
   $78$ linear orders rather than through $411\,840$ six-literal clauses; layers $B$ and $I$ are
   geometric: the boundary polygon of each face is simple, the traversal orientation of each face is
   the one its crossing parity forces, and no vertex carries two reflex corners.
\item \textbf{Conditional non-existence.} If all $59$ formulas are unsatisfiable, the Question has a
   negative answer (Corollary~\ref{cor:conditional}).
\item \textbf{Coverage by cube trees.} A formula is unsatisfiable as soon as every leaf of a finite
   tree of case splits rooted at the empty cube is refuted (Theorem~\ref{thm:coverage}); a refuted
   tree rooted at a nonempty cube settles only that region (Remark~\ref{rem:pilotroot}).
\item \textbf{Verification protocol.} We describe the discipline under which the refutations are to be
   produced and checked, so that the eventual computational claim rests on a formally verified proof
   checker \cite{TanHeuleMyreen2023} rather than on a solver's word, and so that a third party can
   verify it by rechecking retained proofs or, without access to stored proofs, by recomputing any
   or all of its leaves (Section~\ref{sec:protocol}).
\end{enumerate}

\subsection{What is not claimed}\label{sub:notclaimed}
The Boolean computation has \emph{not} been carried out. Theorem~\ref{thm:necessity} and
Corollary~\ref{cor:conditional} are a reduction and a conditional statement; the Question remains
open. Two further limitations are worth stating at the outset. First, only necessity is encoded: a
satisfying assignment of some $\Phi_T$ would not produce a polyhedron, because the layers are
necessary conditions and because a uniform chirotope need not be realizable over $\RR$. Second, the
reduction is stated under the hypotheses (H1)--(H3) of Section~\ref{sec:polyhedra}, which are the
regular polygonal-cell conventions used throughout this paper; dropping them --- allowing faces with holes, two faces to
meet in more than one edge, or additional contacts between faces --- changes the object under study,
and those variants are not settled by these $59$ formulas.

A numerical search for the same twelve-face objects is described by Bokowski and
Kevin~H.\ (CodeParade) \cite[Sections~4 and~8.3]{BokowskiKevin2025}. It found no embedded
realization among the $59$ dual types and presents an intersecting example in Figure~38;
this is a search result, not a non-existence proof. The associated
\textsf{NeighborlyPolyhedra} repository \cite{HackerPoet} documents the search and further
near-misses. Our approach instead gives finite necessary conditions whose certified
unsatisfiability would imply non-existence. The cited 2025 paper leaves the twelve-face
question unresolved.

\subsection{Organization}
Section~\ref{sec:polyhedra} fixes the hypotheses and proves the local facts: trivalence, the Euler
count, and the vertex lemma that at most one of the three corners at a vertex is reflex.
Section~\ref{sec:maps} defines neighbourly maps and realizations and proves the case list.
Section~\ref{sec:genpos} proves the general-position theorem. Section~\ref{sec:dictionary} develops
the sign dictionary. Sections~\ref{sec:chirotope} and \ref{sec:geometric} define and justify the five
layers, Section~\ref{sec:gauge} the gauge. Section~\ref{sec:formula} assembles the formula, states the
reduction and proves that refuted trees of cases cover the formula, and Section~\ref{sec:protocol}
describes the verification protocol.
Section~\ref{sec:status} states precisely what remains to be done. Appendix~\ref{app:numbers} contains
the numerical data, and
Appendix~\ref{app:software} describes the accompanying software.

\subsection*{Acknowledgement of tooling}
The accompanying implementation is written in Python with no runtime dependencies beyond the standard
library; it uses CaDiCaL \cite{BiereEtAl2024} as a proposing solver and \texttt{cake\_lpr} \cite{TanHeuleMyreen2023} as the
accepting checker. The partition heuristic follows \cite{BattlemanReevesHeule2025}.

\section{Polyhedra, hypotheses and local structure}\label{sec:polyhedra}

\begin{definition}\label{def:polyhedron}
A \emph{polyhedron} is a compact closed surface $\Sigma\subset\RR^3$ without self-intersections that
is a finite union of closed planar polygons, its \emph{faces}. Each face is a closed simple polygonal
disk and is the closure in $\Sigma$ of a maximal connected relatively open subset of $\Sigma$
contained in an affine plane. \emph{Edges} are the maximal
straight segments of the union of face boundaries; degree-two subdivision points are suppressed, so
faces and edges form a finite regular cell decomposition of $\Sigma$. Two faces \emph{share an edge}
when their intersection is an edge of both. The polyhedron is \emph{neighbourly} when every two faces
share an edge.
\end{definition}

Three hypotheses are used throughout. Each is a standard convention, and each marks a place where a
wider reading of the word \emph{polyhedron} would leave a case uncovered.

\begin{hypothesis}\label{hyp:H}\leavevmode
\begin{enumerate}[label=\textup{(H\arabic*)},leftmargin=3.2em]
\item The intersection of two faces is a common face of both: empty, a single common vertex, or a
   single common edge. This is our polygonal-cell intersection convention; it forces \emph{exactly} one shared
   edge per pair and, with Lemma~\ref{lem:trivalent}, trivalent vertices.
\item Each face is a simple polygon: one boundary cycle, no holes.
\item Distinct faces have distinct supporting planes.
\end{enumerate}
\end{hypothesis}

\begin{remark}
Within the closed embedded polygonal-cell model of Definition~\ref{def:polyhedron},
maximality of faces is equivalent to requiring that no two edge-adjacent faces be coplanar.
Indeed, two
coplanar faces sharing an edge would have interiors joined through an interior point of that edge
inside a connected relatively open planar subset of $\Sigma$, contradicting maximality.
Conversely, if adjacent faces are noncoplanar, no relatively open planar subset of $\Sigma$
can contain an edge-interior point: its neighbourhood contains open portions of both
incident faces. Nor can such a subset contain a vertex, since any neighbourhood of the
vertex contains edge-interior points. Thus the connected relatively open planar patches
are contained in individual face interiors, whose closures are the faces.
For a neighbourly surface every pair is edge-adjacent, so this condition is precisely (H3).

The model in \cite[Section~1, pp.~1430--1431]{ArsenevaEtAl2024} is broader: its polygonal
complexes may have boundary and may have coplanar cells (indeed, Proposition~3.2 constructs
flat complexes). It imposes neither our maximality convention nor (H3).
The equivalence just proved applies after restricting to a closed embedded surface with
our cell conventions and excluding coplanar adjacent cells; it is not an equivalence with
that broader model. We cite its Section~2, p.~1435, for the open question about closed surfaces.
\end{remark}

\begin{lemma}[Trivalence]\label{lem:trivalent}
In a neighbourly polyhedron satisfying \textup{(H1)}, every vertex has degree three.
\end{lemma}

\begin{proof}
Let $v$ be a vertex. Because $\Sigma$ is a surface with a regular cell decomposition, the faces
incident to $v$ occur in a cyclic order around $v$, consecutive ones sharing an edge through $v$;
a degree-two vertex would make its two incident faces share both incident edges, contrary to (H1).
Thus the degree $d$ of $v$ is at least three. Suppose $d\ge 4$ and let
$F^i,F^j$ be two incident faces that are not consecutive in that cyclic order. Their closed sectors at
$v$ meet only in $v$, so no edge through $v$ is common to both. By neighbourliness they share an edge
$E$, necessarily not incident to $v$. Then $F^i\cap F^j$ contains the edge $E$ and also the point
$v\notin E$. By (H1) it is a single common face of both --- empty, a vertex or an edge --- and,
containing $E$, it would have to be an edge containing $E$, that is $E$ itself, since distinct edges
meet at most in endpoints. But $E$ does not contain $v$, contradicting (H1).
\end{proof}

\begin{proposition}[Euler count]\label{prop:euler}
Let $\Sigma$ be a neighbourly polyhedron with $n\ge 4$ faces satisfying \textup{(H1)}. Then
\[
  E=\binom{n}{2},\qquad V=\frac{n(n-1)}{3},\qquad
  \chi=n-\frac{n(n-1)}{6},\qquad g=\frac{(n-3)(n-4)}{12},
\]
and $g$ is a non-negative integer, so $n\equiv 0,3,4,7\pmod{12}$. For $n=12$: $V=44$, $E=66$,
$\chi=-10$ and $g=6$.
\end{proposition}

\begin{proof}
Each pair of faces shares exactly one edge and each edge lies in exactly two faces, so
$E=\binom{n}{2}$. By Lemma~\ref{lem:trivalent} every vertex has degree three, and counting
vertex--edge incidences gives $3V=2E$, whence $V=n(n-1)/3$. Then $\chi=V-E+F=n-n(n-1)/6$ and
$2-\chi=(n-3)(n-4)/6$. A compact closed surface embedded in $\RR^3$ is orientable --- a closed
non-orientable surface has torsion in $H_1(-;\mathbb{Z})$, which Alexander duality forbids for a
subspace of $\RR^3$; see \cite[\S3.3, Corollary 3.46 and the remarks following it]{Hatcher2002} ---
and it is connected, because any two faces share an edge;
so $\chi=2-2g$ with $g\in\mathbb{Z}_{\ge0}$, and $g=(n-3)(n-4)/12$. Finally $12\mid (n-3)(n-4)$ holds
exactly for $n\equiv0,3,4,7\pmod{12}$, as one checks on the twelve residues.
\end{proof}

The last local fact we need is about reflex corners. Let $v$ be a vertex of a neighbourly polyhedron,
incident to the three faces $F^f,F^g,F^h$. Near $v$ each of them is a plane wedge, and the corner of
$F^f$ at $v$ is \emph{reflex} when its interior angle exceeds $\pi$.

\begin{lemma}[At most one reflex corner]\label{lem:onereflex}
At every vertex of degree three of a polyhedron satisfying \textup{(H1)}--\textup{(H3)}, at most one
of the three corners is reflex.
\end{lemma}

\begin{proof}
Let $v$ be incident to $F^f,F^g,F^h$ and suppose the corners of $F^f$ and $F^g$ at $v$ are both
reflex. Both faces contain the edge $E(f,g)$, which lies on the line $L=h_f\cap h_g$ and has $v$ as an
endpoint; let $R$ be the ray of $L$ issuing from $v$ opposite to $E(f,g)$. Note first that
$F^f\cap F^g=E(f,g)$: by (H1) that intersection is a common face of both --- empty, a vertex or an
edge --- and it contains the edge $E(f,g)$.

Near $v$ the face $F^f$ is the wedge bounded by the germs of $E(f,g)$ and $E(f,h)$. A plane wedge has
interior angle exceeding $\pi$ if and only if its interior contains the ray opposite to each of its
two bounding rays; in particular a reflex corner of $F^f$ at $v$ puts a germ of $R$ in the interior of
$F^f$. The same applies to $F^g$. Hence $F^f\cap F^g$ contains points of $R$ other than $v$, while $E(f,g)$ meets $R$ only in $v$; a
contradiction.
\end{proof}

\begin{remark}\label{rem:octants}
The software additionally checks the local conclusion by enumerating the $256$ unions of
the eight octants cut out by three independent planes. The admissible sector-size profiles
are $(1,1,1)$ and $(1,1,3)$, occurring $16$ and $48$ times respectively; an entry $3$
means a reflex corner. The enumeration and its incidence conventions are documented with
the code. It is a supplementary check; the reduction uses the proof above.
\end{remark}

\section{Maps, realizations and the case list}\label{sec:maps}

\subsection{Neighbourly maps}

\begin{definition}\label{def:map}
A \emph{neighbourly map} on $n\ge4$ faces is a set $T$ of $3$-element subsets of $\{1,\dots,n\}$,
called \emph{vertices}, such that
\begin{enumerate}[label=\textup{(M\arabic*)},leftmargin=3.2em]
\item every pair $\{i,j\}$ lies in exactly two members of $T$, written $\Ends(i,j)=\{k,l\}$ where
   $\{i,j,k\},\{i,j,l\}\in T$; and
\item for every $f$ the \emph{link} $\link(f)$ --- the graph on $\{1,\dots,n\}\setminus\{f\}$ with
   $a\sim b$ iff $\{f,a,b\}\in T$ --- is a single cycle of length $n-1$.
\end{enumerate}
The map is \emph{orientable} when the $|T|=n(n-1)/3$ triples can be cyclically oriented so that every
ordered pair $(a,b)$ with $a\ne b$ occurs in exactly one oriented triple.
\end{definition}

Condition (M1) says $|T|=n(n-1)/3$ and makes $T$ the triangle set of a two-dimensional simplicial
complex on the labels $1,\dots,n$ in which every two labels span an edge and every edge lies in
exactly two triangles. Each link $\link(f)$ is then $2$-regular, a disjoint union of cycles, and a
label whose link has more than one cycle would be a singular point of the complex. Condition (M2)
excludes such points: it makes the complex a triangulation of a closed surface. The following lemma
is what lets us pass from geometry to $T$.

\begin{lemma}\label{lem:inducedmap}
Let $\Sigma$ be a neighbourly polyhedron with $n$ faces satisfying \textup{(H1)}--\textup{(H2)}, its
faces labelled $1,\dots,n$. Then the set $T$ of triples of faces meeting at a vertex is a neighbourly
map on $n$ faces, and distinct vertices give distinct triples.
\end{lemma}

\begin{proof}
By Lemma~\ref{lem:trivalent} each vertex is incident to exactly three faces, so each vertex gives a
$3$-subset. Each pair $\{i,j\}$ shares exactly one edge, and that edge has exactly two endpoints, so
$\{i,j\}$ lies in exactly two triples counted with multiplicity: (M1) holds once we know the triples
of distinct vertices are distinct. Suppose $v\ne w$ were vertices with the same incident faces
$\{f,g,h\}$. Then $v$ and $w$ are both endpoints of the edge $E(f,g)$, and likewise of $E(f,h)$ and
of $E(g,h)$. Three distinct straight segments cannot have the same pair of endpoints, so two of these
edges coincide, contradicting the regularity of the cell decomposition. Finally (M2) is the statement
that the faces around a fixed face $f$ close up into one cycle, which holds because the boundary of
the disk $F^f$ is a single closed polygon whose edges are $E(f,u)$, $u\ne f$, each consecutive pair
meeting at a vertex.
\end{proof}

\subsection{Realizations}
Write a plane as its covector: $h=\{x\in\RR^3: q\cdot(x,1)=0\}$ with
$q=(a,b,c,d)\in\RR^4$, $(a,b,c)\ne 0$; its \emph{normal} is $n(q)=(a,b,c)$. Adjoin the
\emph{plane at infinity} $q_0=(0,0,0,1)$, which no affine point satisfies. For a $4$-subset
$S\subseteq\{0,1,\dots,n\}$ let
\[
  \bra{S} \;=\; \dett\big(q_s\big)_{s\in S}
\]
with the rows taken in increasing order of the labels; for an ordered tuple the bracket changes sign
with the parity of the permutation. Write $P_{ijk}=h_i\cap h_j\cap h_k$ when that intersection is a
single point and $L_{ij}=h_i\cap h_j$ when that intersection is a line.

\begin{definition}\label{def:realization}
A \emph{realization} of a neighbourly map $T$ on $n$ faces is a tuple $q=(q_1,\dots,q_n)$ of covectors
such that, writing $E(i,j)$ for the segment of $L_{ij}$ between $P_{ij k}$ and $P_{ij l}$ with
$\{k,l\}=\Ends(i,j)$, and $F^f$ for the closed region bounded in $h_f$ by the closed polygon with
vertex sequence $\big(P_{f,u,\,\link\text{-successor of }u}\big)_{u\in\link(f)}$:
\begin{enumerate}[label=\textup{(R\arabic*)},leftmargin=3.2em]
\item for every $\{i,j,k\}\in T$ the normals $n(q_i),n(q_j),n(q_k)$ are linearly independent, so
  $P_{ijk}$ exists, and the $n(n-1)/3$ points so obtained are pairwise distinct;
\item for every $f$ the polygon above is simple, so $F^f$ is a closed disk;
\item $F^i\cap F^j=E(i,j)$ for all $i\ne j$; and
\item the planes $h_1,\dots,h_n$ are pairwise distinct.
\end{enumerate}
\end{definition}

\begin{lemma}\label{lem:bridge}
The union $\Sigma=\bigcup_f F^f$ of the faces of a realization of a neighbourly map $T$ is a
neighbourly polyhedron with $n$ faces satisfying \textup{(H1)}--\textup{(H3)}, whose induced map is
$T$. Conversely every such polyhedron, with its faces labelled, arises this way from the tuple of its
supporting covectors.
\end{lemma}

\begin{proof}
Let $q$ be a realization. Each $F^f$ is a closed disk by (R2), and two of them meet exactly in the
common edge $E(i,j)$ by (R3), which is an edge of both by construction.
No fourth face $F^i$ contains a vertex $P_{fgh}$: by (R3) this point would lie on the three distinct
boundary edges $E(i,f),E(i,g),E(i,h)$ of the simple polygon $F^i$, whereas a point of a simple polygon
belongs to at most two of its edges. At $P_{fgh}$ the three incident edge lines are distinct by
(R1), so no two incident segments join into one straight edge; consequently the edges are exactly
the segments $E(i,j)$. Thus (H1) holds and the faces are pairwise adjacent.
Nor can a third face $F^k$ meet the relative interior of $E(i,j)$: by (R3) the point would lie
on both $E(k,i)$ and $E(k,j)$, which by (R2) can meet only at their common polygon vertex
$P_{ijk}$, an endpoint of $E(i,j)$.
(H2) is (R2) and (H3) is (R4). $\Sigma$ is compact and closed as a finite union
of disks whose boundary edges are each shared by exactly two of them. It is a $2$-manifold: an
interior point of a face has a disk neighbourhood; a point in the relative interior of $E(i,j)$ has a
neighbourhood consisting of two half-disks glued along a segment; and at a vertex $P_{fgh}$ the three
wedges of $F^f,F^g,F^h$ meet pairwise in the germs of $E(f,g),E(f,h),E(g,h)$, so the link of the
vertex in $\Sigma$ is a triangle, that is, a circle. The relative interior of each $F^f$ is a
connected relatively open planar subset of $\Sigma$. Every relatively open neighbourhood in
$\Sigma$ of a point of $\partial F^f$ contains two-dimensional pieces of two incident faces with
distinct supporting planes, by (R4), and therefore is not planar. Thus this interior is maximal
among connected relatively open planar subsets of $\Sigma$, and its closure is $F^f$.
Conversely, a neighbourly polyhedron with labelled faces gives covectors $q_i$ of its supporting
planes, and Lemma~\ref{lem:inducedmap} gives the map. At an incident triple,
the three normals are linearly independent. Indeed, otherwise their rank is two:
the planes are distinct and pass through the same vertex. Their pairwise intersection
lines would then coincide. The three distinct incident edge germs would lie on the
two rays of that line, forcing two germs to overlap, contrary to the regular embedded
cell structure. Thus each incident triple meets at a unique point, and the vertex of $\Sigma$ with
incident faces $\{i,j,k\}$ is $P_{ijk}$; by Lemma~\ref{lem:inducedmap} these $n(n-1)/3$ vertices are
distinct, which is (R1). The common edge of $F^i$ and $F^j$ lies on $L_{ij}$ and its endpoints are the
two vertices whose triples contain $\{i,j\}$, so it is the segment $E(i,j)$ of
Definition~\ref{def:realization}. As in the proof of Lemma~\ref{lem:inducedmap}, the boundary of the
face $F^f$ consists of the edges $E(f,u)$, $u\ne f$, consecutive ones meeting at the vertices
$P_{f,u,\mathrm{succ}(u)}$ in the order of $\link(f)$; it is therefore the polygon of
Definition~\ref{def:realization}, which is simple by (H2), and the face, a closed disk in $h_f$, is the
region that polygon bounds. This is (R2), and the regions $F^f$ of the definition are the faces of
$\Sigma$. By (H1) the intersection $F^i\cap F^j$ is a single common face of both; it contains the edge
$E(i,j)$, and distinct edges meet at most in endpoints, so $F^i\cap F^j=E(i,j)$, which is (R3).
Finally (R4) is (H3).
\end{proof}

\subsection{The case list}

\begin{lemma}[Oriented form]\label{lem:oriented}
Let $\vec T$ be a set of $n(n-1)/3$ cyclically ordered triples containing every ordered pair $(a,b)$,
$a\ne b$, in exactly one triple. Then the underlying sets of $\vec T$ satisfy \textup{(M1)} and
\textup{(M2)} if and only if for every $v$ the map $\sigma_v\colon b\mapsto c$, where
$(v,b,c)\in\vec T$, is an $(n-1)$-cycle on $\{1,\dots,n\}\setminus\{v\}$. The orientable
neighbourly maps on $n$ faces are exactly the sets of underlying sets of such $\vec T$.
\end{lemma}

\begin{proof}
In both directions, $\sigma_v$ is a permutation of the $n-1$ elements other than $v$, because each
ordered pair $(v,b)$ lies in exactly one member of $\vec T$.

Suppose every $\sigma_v$ is an $(n-1)$-cycle. Then no $\sigma_v$ has a $2$-cycle, so no triple occurs
in $\vec T$ with both of its cyclic orders: from $(a,b,c)$ and $(a,c,b)$ we would get $\sigma_a(b)=c$
and $\sigma_a(c)=b$. Hence the $n(n-1)/3$ underlying sets are distinct, and because the two ordered
versions of an unordered pair lie in two different members of $\vec T$, every unordered pair lies in
exactly two of those sets: that is (M1). The cycles of $\sigma_v$ are the connected components of
$\link(v)$, so (M2) holds as well.

Conversely, suppose the underlying sets satisfy (M1) and (M2). Counting pair--triple incidences, (M1)
forces exactly $n(n-1)/3$ underlying sets, so no triple occurs in $\vec T$ twice and no $\sigma_v$
has a $2$-cycle; the cycles of $\sigma_v$ are then the components of $\link(v)$, and (M2) says
precisely that $\sigma_v$ is an $(n-1)$-cycle.

Finally, a coherent orientation of the triangles of a triangulated closed surface induces on each edge
the two opposite directions from its two incident triangles, which is exactly the condition that every
ordered pair occurs once; conversely such a set of oriented triples orients the dual surface. So the
orientable neighbourly maps are exactly those obtained this way.
\end{proof}

\begin{lemma}[Framing count]\label{lem:framing}
Let $\mathcal{N}_n$ be the set of \emph{orientable} neighbourly maps on the label set
$\{0,1,\dots,n-1\}$ whose link at $0$ is exactly the cycle $(1,2,\dots,n-1)$. Then
\[
  |\mathcal{N}_n| \;=\; \sum_{C}\frac{2n(n-1)}{|\Aut(C)|},
\]
the sum being over isomorphism classes $C$ of orientable neighbourly maps on $n$ faces. Moreover any isomorphism
between two members of $\mathcal{N}_n$ is determined by the image of the label $0$ together with the
rotation or reflection it induces on the link of $0$, so at most $2n(n-1)$ relabellings have to be
tried to decide isomorphism.
\end{lemma}

\begin{proof}
Orientability is invariant under relabelling, so it is enough to argue inside one class. Fix a class
$C$ with representative $M_0$ and $a=|\Aut(C)|$. The labelled maps in $C$ are the
$M_0^{\sigma}$, $\sigma\in S_n$, and $M_0^{\sigma}=M_0^{\tau}$ exactly when $\sigma\tau^{-1}\in\Aut$.
Because the map is neighbourly, the link of a vertex $v$ of $M_0$ consists of all other $n-1$
vertices, so a relabelling $\sigma$ with $M_0^\sigma\in\mathcal N_n$ is determined by the vertex
$v=\sigma^{-1}(0)$ --- $n$ choices --- together with a bijection from $\link(v)$ to
$\{1,\dots,n-1\}$ carrying the link cycle of $v$ to $(1,2,\dots,n-1)$, of which there are exactly
$2(n-1)$, one for each starting point and direction. Hence $2n(n-1)$ relabellings, and the members of
$\mathcal N_n$ they produce are counted $a$ times each. The last claim is the same computation read
backwards.
\end{proof}

\begin{remark}
Non-orientable neighbourly maps on twelve faces also exist:
\cite[Theorem~4]{EllinghamStephens2005} enumerates triangular embeddings of $K_{12}$
on the non-orientable surface of the same Euler characteristic.
This is why orientability is part of the definition of $\mathcal{N}_n$. They are
irrelevant here, because the surface of a polyhedron is orientable
(Proposition~\ref{prop:euler}).
\end{remark}

\begin{theorem}[Case list]\label{thm:59}
Every neighbourly polyhedron with twelve faces satisfying \textup{(H1)}--\textup{(H3)} induces one of
exactly $59$ pairwise non-isomorphic orientable neighbourly maps on twelve faces.
\end{theorem}

\begin{proof}
By Lemma~\ref{lem:inducedmap} the polyhedron induces a neighbourly map $T$.
Its dual is a triangulation with twelve vertices and complete edge graph $K_{12}$.
It is orientable of genus $6$ by Proposition~\ref{prop:euler}, since it carries the
dual cell structure of the same surface. The classification of
\cite{AltshulerBokowskiSchuchert1996}, also reproduced in
\cite[Theorem 4]{EllinghamStephens2005}, gives exactly $59$ isomorphism classes
of these orientable triangulations. Taking their combinatorial duals gives the claim.
\end{proof}

\paragraph{Recomputing the representatives.}
The implementation enumerates these maps without reading an external list.
Lemma~\ref{lem:oriented} reduces the search to cyclically oriented triples covering
each ordered pair once. Fix the oriented link of one label to $(1,2,\dots,11)$,
as allowed by Lemma~\ref{lem:framing}; this fixes eleven of the forty-four triples.
At each step choose an uncovered ordered pair $(a,b)$ and try every third label $c$
for which $(b,c)$ and $(c,a)$ are also uncovered. Reject an addition if it closes
any partial link permutation into a cycle shorter than $11$.

These rejections lose no completion: an ordered pair cannot be covered twice,
and a closed proper link cycle cannot later be extended to a single $11$-cycle.
Conversely, every completed branch covers every ordered pair and has a single
cycle at every label, so Lemma~\ref{lem:oriented} makes it an orientable map.
Every target map has a labelling with the fixed link and a unique coherent
orientation extending that link; at each chosen pair its next triple is among
the candidates. Induction therefore shows that the search visits every such
map. At most forty-four triples are placed, so the search is finite.
Canonicalization tries the $264$ relabellings of Lemma~\ref{lem:framing}.
The implemented search returns $12\,144$ framed maps in $59$ classes and checks
that their canonical forms are exactly those of the shipped matrices.
Appendix~\ref{app:numbers} gives the automorphism counts as reproducibility checks.

\paragraph{Identification with the published list.}
The formulas use the $59$ matrices shipped with this project, in the project's
own numbering. Their duals are also exactly the $59$ triangulations listed by
Lutz \cite{LutzManifoldPage}. For each entry of that list, the accompanying data
give a bijection of the twelve labels taking one shipped map to that entry;
every map is used once. Applying these bijections reproduces the published facet
lists, whose normalized SHA-256 is recorded with the comparison. This check
uses the explicit bijections, not the enumeration's canonicalization algorithm.
The Manifold Page cites \cite{SulankeLutz2009,EllinghamStephens2005} for its
twelve-vertex enumerations. No external list is needed to generate the formulas
or to rerun the census.

Thus the classification supplies the completeness theorem, the search supplies
a reproducible recomputation, and the explicit correspondence identifies the
concrete input data with a published list. Agreement of counts alone would not
identify the maps. The orbit identity and the small-size controls test the
implementation; they are not the completeness argument. The classification
itself is an existing computational classification, not a new classification
proved by the SAT computation proposed here.

\section{General position}\label{sec:genpos}

Everything that follows reads a realization through the signs of its brackets, which requires those
signs to be non-zero.

\begin{hypothesis}[General position]\label{hyp:G}\leavevmode
\begin{enumerate}[label=\textup{(G)},leftmargin=3.2em]
\item $\bra{S}\ne0$ for every $4$-subset $S\subseteq\{0,1,\dots,n\}$.
\end{enumerate}
\end{hypothesis}

Geometrically: no four of the $n+1$ covectors are linearly dependent. The brackets containing $0$ say
that any three of the planes $h_1,\dots,h_n$ meet in exactly one affine point, because
$\bra{0,i,j,k}=-\dett(n_i,n_j,n_k)$ by Lemma~\ref{lem:cofactor}; the brackets not containing $0$ say
that no four of the planes share a point of projective space --- no four pass through a common affine
point, and no four are parallel to a common direction. For $n=12$ that is $220$ conditions of the
first kind and $495$ of the second.

\begin{theorem}[Reduction to general position]\label{thm:genpos}
Every neighbourly map admitting a realization satisfying \textup{(R1)}--\textup{(R4)} also admits
one satisfying \textup{(G)}.
\end{theorem}

The main point is to preserve the embedded surface when its supporting planes move. We use an
auxiliary triangulation and the following elementary stability fact; compare
\cite[Section 2, Lemmas 2.1--2.2]{Lazarus2020}. The triangulation is used only in this proof and
does not change the faces of the map or the Boolean encoding.

\begin{lemma}[Stability of a finite triangulation]\label{lem:plstability}
Let $K$ be a finite geometric simplicial complex in $\RR^d$. Every sufficiently small displacement
of its vertices extends, affinely on each simplex, to an embedding $F\colon |K|\to\RR^d$.
\end{lemma}

\begin{proof}
First, such a map is injective if the images of every two simplices with disjoint vertex sets
are disjoint. To see this, suppose $F(x)=F(y)$ for distinct $x,y\in |K|$. Write their barycentric
coordinates as $(\alpha_v)$ and $(\beta_v)$ over the vertices of $K$, with zero coefficients
outside their supporting simplices. The coefficients $c_v=\alpha_v-\beta_v$ are not all zero
and satisfy
\[
  \sum_v c_v=0,\qquad \sum_v c_v F(v)=0.
\]
The positive coefficients are supported on a face of the simplex supporting $x$, and the
negative coefficients on a face of the simplex supporting $y$. These two faces have disjoint
vertex sets. Dividing by $\sum_{c_v>0}c_v=-\sum_{c_v<0}c_v>0$ gives a point in both images,
a contradiction.

In the original geometric complex, any two simplices with disjoint vertex sets are disjoint
compact sets, hence at positive distance. If there are such pairs, let $\delta>0$ be the minimum
of their finitely many distances. Move every vertex by less than $\delta/3$. A point of any
simplex, expressed as a convex combination of its vertices, moves by less than $\delta/3$;
therefore the images of each such pair remain disjoint. The criterion above proves injectivity.
If there are no such pairs, the criterion applies without a distance bound. The affine pieces
agree on common faces, so $F$ is continuous; compactness of $|K|$ makes it an embedding.
\end{proof}

\begin{proof}[Proof of Theorem~\ref{thm:genpos}]
Fix a realization $q^*=(q_1^*,\dots,q_n^*)$. Triangulate each of its simple polygonal faces by
diagonals, using its existing vertices and nondegenerate triangles. Such a triangulation exists
also for nonconvex simple polygons \cite{Meisters1975}. By (R3), the triangulations fit together
along the original edges to form a finite geometric simplicial complex $K$: two faces $i\ne j$ meet
only in the segment $E(i,j)$, an edge of both polygons that no triangulation subdivides, and a
triangle of either face meets that segment in the empty set, in one endpoint, or in the whole
segment as one of its edges; so a triangle of face $i$ and a triangle of face $j$ meet in a common
face. Let $K_f$ be the subcomplex triangulating face $f$.

For $q$ close to $q^*$, move each original vertex $v=\{i,j,k\}\in T$ to
\[
  p_v(q)=P_{ijk}(q)=h_i(q)\cap h_j(q)\cap h_k(q).
\]
These points exist and depend continuously on $q$: the incident normal determinants are nonzero
at $q^*$ by (R1), and Cramer's rule applies in a neighbourhood. There are finitely many vertices,
so Lemma~\ref{lem:plstability} shows that the resulting affine map $F_q$ on $K$ is an embedding
for all sufficiently close $q$.

Every vertex of $K_f$ lies in $h_f(q)$ after the displacement. Thus $F_q(|K_f|)$ is an embedded
planar disk with the prescribed boundary polygon, and equals the closed region bounded by that
polygon. Indeed, $F_q$ maps the boundary circle of $|K_f|$ injectively onto that polygon, which is
therefore simple, and by invariance of domain it maps the interior of $|K_f|$ onto an open subset
of the plane $h_f(q)$; this set is connected, bounded and closed in the complement of the
polygon, so it is the bounded component of that complement. In particular, (R2) holds.
Injectivity gives, for $i\ne j$,
\[
  F_q(|K_i|)\cap F_q(|K_j|)
  =F_q\bigl(|K_i|\cap |K_j|\bigr)=E(i,j)(q),
\]
which is (R3). The incident normal determinants remain nonzero, and distinct vertices remain
distinct by injectivity, giving (R1). The normals remain nonzero and the supporting planes
remain pairwise distinct by continuity, giving (R4). Hence a whole neighbourhood of $q^*$ in
$Q=(\RR^4)^n$ consists of realizations of the same map.

For each $4$-subset $S\subseteq\{0,1,\dots,n\}$, the bracket $\bra{S}$ is a nonzero polynomial
on $Q$, with $q_0=(0,0,0,1)$ fixed. Indeed, its value can be made nonzero by choosing the four
rows to be a basis of $\RR^4$; when $0\in S$, choose the other rows to be $e_1,e_2,e_3$.
The product of these finitely many polynomials is nonzero, so its zero set contains no open ball.
The neighbourhood just obtained therefore contains a realization at which all brackets are
nonzero, as required.
\end{proof}

\begin{remark}
The argument also shows that the new realization can be chosen arbitrarily close to the original
one. It moves the supporting planes and derives the vertices from their intersections; arbitrary
independent vertex displacements would not preserve the planarity of the original faces.
\end{remark}

We record two local geometric facts used in the sign dictionary below; in them $q^{*}$ is any
realization of $T$.

\begin{lemma}[Nondegeneracy is open]\label{lem:step1}
At $q^{*}$ the normals $n(q_i),n(q_j),n(q_k)$ are independent for every $\{i,j,k\}\in T$ and
$n(q_i),n(q_j)$ are independent for every $i\ne j$. Consequently, on a neighbourhood of $q^{*}$ the
vertex maps $q\mapsto P_{ijk}(q)$ and the line maps $q\mapsto L_{ij}(q)$ are real-analytic, and all
vertices stay inside one fixed ball.
\end{lemma}

\begin{proof}
Independence at a vertex is (R1). For a pair, if $n(q_i)$ and $n(q_j)$ were parallel the planes
would be equal, excluded by (R4), or disjoint, and then $E(i,j)\subseteq h_i\cap h_j$ could not be a
segment. Both conditions are the non-vanishing of a determinant, hence open, and on the set where they
hold Cramer's rule expresses $P_{ijk}$ and a direction vector and base point of $L_{ij}$ as rational
functions of $q$ with non-vanishing denominators.
\end{proof}

\begin{lemma}[No straight angles]\label{lem:step2}
At $q^{*}$ no three consecutive vertices of a face polygon are collinear. Hence on a neighbourhood the
turn at every corner is non-zero and of constant sign.
\end{lemma}

\begin{proof}
Three consecutive vertices of the polygon of $f$ are $P_{f,a,x}$, $P_{f,a,b}$, $P_{f,b,y}$ with $a,b$
consecutive in $\link(f)$. The first two lie on $L_{fa}$ and are distinct by (R1); the last two lie on
$L_{fb}$ and are distinct. If the three were collinear, the line through them would equal both
$L_{fa}$ and $L_{fb}$, so $h_f,h_a,h_b$ would share a line and their normals would be dependent,
contradicting Lemma~\ref{lem:step1}. The turn at a corner is a continuous function of $q$ which is
non-zero at $q^{*}$, hence of constant sign nearby.
\end{proof}

\begin{convention}\label{conv:genpos}
From now on all realizations are assumed to satisfy \textup{(G)}. By Theorem~\ref{thm:genpos} this
loses no map: if some map is realizable at all, it is realizable in general position.
\end{convention}

\section{The sign dictionary}\label{sec:dictionary}

Throughout this section $q$ is a realization of a neighbourly map $T$ on $n$ faces satisfying (G). For
a covector $q_u$ write $h_u(x)=q_u\cdot(x,1)$ for the affine form whose zero set is the plane, and
$n_u=n(q_u)$ for its normal. For $S$ a $4$-subset of $\{0,1,\dots,n\}$ put
\[
  X_S=\begin{cases}1,&\bra{S}<0,\\ 0,&\bra{S}>0,\end{cases}
\]
which is defined by (G); for an ordered tuple, $X$ changes by the parity of the permutation, so
$X_{\sigma(S)}=X_S\xor\operatorname{par}(\sigma)$. All the identities below are exact, and all of them
express a geometric predicate as an $\FF$-linear or small Boolean function of the $\binom{n+1}{4}$
bits $X_S$.

\subsection{Two cofactor identities}

\begin{lemma}\label{lem:cofactor}
For distinct $i,j,k,l\in\{1,\dots,n\}$ with $n_i,n_j,n_k$ independent,
\[
  \bra{0,i,j,k}=-\dett(n_i,n_j,n_k),
  \qquad
  \bra{i,j,k,l}=\dett(n_i,n_j,n_k)\,h_l(P_{ijk}).
\]
\end{lemma}

\begin{proof}
The first is the expansion of the determinant along the row $q_0=(0,0,0,1)$, whose only non-zero entry
sits in the last column with sign $(-1)^{1+4}$. For the second, change affine coordinates by
$y=x-P_{ijk}$. A covector $q_u=(n_u,d_u)$ becomes $(n_u,h_u(P_{ijk}))$, because
$q_u\cdot(x,1)=n_u\cdot y+h_u(P_{ijk})$; this is right multiplication of the row by a unipotent matrix,
so the $4\times4$ determinant is unchanged. In the new coordinates the first three rows have last
entry $0$, and expanding along the last column leaves
$h_l(P_{ijk})\cdot\dett(n_i,n_j,n_k)$.
\end{proof}

\subsection{Order along a line}
Fix $i\ne j$ and orient the line $L_{ij}$ by the direction $d_{ij}=n_i\times n_j$, which is non-zero by
Lemma~\ref{lem:step1}. The other $n-2$ planes mark $n-2$ points $P_{ijk}$ on it, pairwise distinct by
(G). Let $\Ord(i,j;k,l)=1$ when $P_{ijk}$ precedes $P_{ijl}$ in that direction, and $0$ otherwise.

\begin{theorem}[Order bits]\label{thm:order}
For all distinct $i,j,k,l\in\{1,\dots,n\}$,
\[
   \Ord(i,j;k,l)\;=\;X_{0ijk}\xor X_{0ijl}\xor X_{ijkl}.
\]
\end{theorem}

\begin{proof}
Write $D_{ijk}=\dett(n_i,n_j,n_k)$ and $P_{ijl}-P_{ijk}=t\,d_{ij}$; then $\Ord(i,j;k,l)=[t>0]$. Since
$h_l(P_{ijl})=0$ and $n_l\cdot d_{ij}=\dett(n_l,n_i,n_j)=D_{ijl}$,
\[
  h_l(P_{ijk})=h_l\big(P_{ijl}-t\,d_{ij}\big)=-t\,D_{ijl},
\]
so Lemma~\ref{lem:cofactor} gives $\bra{i,j,k,l}=-t\,D_{ijk}D_{ijl}$ and
$\bra{0,i,j,k}=-D_{ijk}$, $\bra{0,i,j,l}=-D_{ijl}$. All factors are non-zero by (G). Writing
$\alpha=[D_{ijk}<0]$, $\beta=[D_{ijl}<0]$ and $\gamma=[t<0]$ we get $X_{0ijk}=1\xor\alpha$,
$X_{0ijl}=1\xor\beta$ and $X_{ijkl}=(1\xor\gamma)\xor\alpha\xor\beta$, because the sign bit of a
product is the sum of the sign bits. Adding the three gives $1\xor\gamma=[t>0]$.
\end{proof}

\begin{remark}[Why a $0$ appears]\label{rem:cut}
A projective line carries no linear order; what a rank-$4$ chirotope gives on $L_{ij}$ is a
\emph{circular} order of the marked points. Removing one point turns it into a linear order, and the
point removed in Theorem~\ref{thm:order} is the line's own point at infinity, $P_{0ij}$. This is why
the element $0$ occurs in the formula, and recognizing it is what makes Layer $A$ of
Section~\ref{sec:chirotope} possible.
\end{remark}

\subsection{Betweenness inside an edge}

\begin{definition}\label{def:onb}
For $i\ne j$ with $\Ends(i,j)=\{z_0,z_1\}$ and $u\notin\{i,j,z_0,z_1\}$ put
\[
   \OnB(i,j;u)\;=\;\Ord(i,j;z_0,u)\xor\Ord(i,j;z_1,u).
\]
\end{definition}

\begin{lemma}\label{lem:onb}
$\OnB(i,j;u)=1$ if and only if $P_{iju}$ lies in the relative interior of the edge $E(i,j)$.
\end{lemma}

\begin{proof}
$E(i,j)$ is the segment of $L_{ij}$ between $P_{ijz_0}$ and $P_{ijz_1}$. In a strict total order,
exactly one of the two endpoints precedes a given point precisely when that point lies strictly
between them; by (G) the point $P_{iju}$ coincides with neither endpoint, since that would make four
planes concurrent.
\end{proof}

\begin{lemma}[Cyclic adjacency]\label{lem:adjacency}
Two edges $E(f,a)$ and $E(f,b)$ of the face $f$ are consecutive in the boundary polygon of $F^f$ if
and only if $\{f,a,b\}\in T$, and then they meet exactly at the vertex $P_{fab}$.
\end{lemma}

\begin{proof}
By construction the boundary polygon of $F^f$ visits the faces of $\link(f)$ in the cyclic order of
that link, and $a,b$ are adjacent in $\link(f)$ exactly when $\{f,a,b\}\in T$
(Definition~\ref{def:map}). Consecutive edges meet at the corresponding vertex, and only there by
simplicity of the polygon.
\end{proof}

\subsection{Turns, orientation and reflex corners}
Let $a,b$ be consecutive in $\link(f)$, in the direction of a fixed traversal of the link cycle, and
let $v=P_{fab}$ be the corner between the edges $E(f,a)$ and $E(f,b)$. Let $x$ be the third face at the
other endpoint of $E(f,a)$ and $y$ the third face at the other endpoint of $E(f,b)$, that is
$\Ends(f,a)=\{b,x\}$ and $\Ends(f,b)=\{a,y\}$. Put
\[
   u_{\mathrm{in}}=v-P_{fax},\qquad u_{\mathrm{out}}=P_{fby}-v,\qquad
   w(f;a,b)=\big(u_{\mathrm{in}}\times u_{\mathrm{out}}\big)\cdot n_f .
\]
By Lemma~\ref{lem:step2}, $w(f;a,b)\ne0$.

\begin{lemma}[Turn bits]\label{lem:turn}
$\big[w(f;a,b)<0\big]\;=\;1\xor\Ord(f,a;x,b)\xor\Ord(f,b;a,y)\xor X_{0fab}$.
\end{lemma}

\begin{proof}
Orient $L_{fa}$ by $d_{fa}=n_f\times n_a$, the direction used in Theorem~\ref{thm:order}. Then
$u_{\mathrm{in}}=s\,d_{fa}$ with $[s<0]=1\xor\Ord(f,a;x,b)$, because $\Ord(f,a;x,b)=[P_{fax}<P_{fab}]$
and $u_{\mathrm{in}}=P_{fab}-P_{fax}$. Likewise $u_{\mathrm{out}}=t\,d_{fb}$ with
$[t<0]=1\xor\Ord(f,b;a,y)$. From the identity $(A\times B)\times(A\times C)=\dett(A,B,C)\,A$,
\[
  \big(d_{fa}\times d_{fb}\big)\cdot n_f=\big((n_f\times n_a)\times(n_f\times n_b)\big)\cdot n_f
  =\dett(n_f,n_a,n_b)\,\lVert n_f\rVert^{2},
\]
and $\dett(n_f,n_a,n_b)=-\bra{0,f,a,b}$ by Lemma~\ref{lem:cofactor}, so
$[\dett(n_f,n_a,n_b)<0]=1\xor X_{0fab}$. Since $w=s\,t\,(d_{fa}\times d_{fb})\cdot n_f$ and the sign
bit of a product is the sum of the sign bits,
\[
 [w<0]=(1\xor\Ord(f,a;x,b))\xor(1\xor\Ord(f,b;a,y))\xor(1\xor X_{0fab}),
\]
which is the claim.
\end{proof}

\begin{lemma}[One orientation bit per face]\label{lem:sf}
For each face $f$ there is $s_f\in\FF$, depending on $f$ and the chosen traversal only, such that for
every corner $(a,b)$ of $f$
\[
   \refl(f;a,b)\;:=\;\big[w(f;a,b)<0\big]\xor s_f
\]
equals $1$ exactly when the interior angle of $F^f$ at $P_{fab}$ exceeds $\pi$. We write
$\refl(f;a,b)$ also when the pair is given in the opposite order, meaning the same bit: the pair is
normalized to link-cycle order before the formula of Lemma~\ref{lem:turn} is applied.
\end{lemma}

\begin{proof}
Orient the plane $h_f$ by $n_f$. A simple closed polygon in an oriented plane is traversed either
positively or negatively; let $s_f=0$ in the first case and $s_f=1$ in the second. For a positively
traversed simple polygon the interior angle at a corner exceeds $\pi$ exactly when the turn from the
incoming to the outgoing direction is negative, that is when $w<0$; for a negatively traversed one the
condition is $w>0$. Both cases are summarized by $[w<0]\xor s_f$. The bit $s_f$ does not depend on the
corner, because the traversal direction is a property of the polygon.
\end{proof}

\begin{remark}[An easy way to get this wrong]\label{rem:order-matters}
Both sides of Lemma~\ref{lem:turn} change under exchanging $a$ and $b$: the corner is the same, but the
formula computes the complement. The predicate ``this corner is reflex'' is symmetric while its
expression is antisymmetric, and the two are reconciled only by always writing the pair in the
direction of the link traversal. An implementation that reaches a corner in the other order produces a
well-formed formula that encodes something else.
\end{remark}

\begin{lemma}[Crossing parity fixes the orientation]\label{lem:parity}
Let $f$ be a face, $g\in\link(f)$, $\Ends(f,g)=\{z_0,z_1\}$ and $V=P_{f,g,z_1}$. For
$v\notin\{f,g,z_0,z_1\}$ put
\[
   C_v=\OnB(f,v;g),\qquad B_v=\Ord(f,g;z_1,z_0)\xor\Ord(f,g;z_1,v).
\]
Then the corner of $F^f$ at $V$ is reflex if and only if $\bigoplus_{v} C_v\wedge B_v=1$.
\end{lemma}

\begin{proof}
Let $R$ be the open ray of $L_{fg}$ issuing from $V$ in the direction opposite to the edge $E(f,g)$;
thus $B_v=1$ says exactly that the point $P_{fgv}$, which is where the line $L_{fv}$ meets $L_{fg}$,
lies on $R$: indeed $B_v$ is the exclusive or of $[V<P_{fgz_0}]$ and $[V<P_{fgv}]$, which says that
$P_{fgv}$ and the other endpoint $P_{fgz_0}$ of $E(f,g)$ lie on opposite sides of $V$.

A plane wedge contains the ray opposite to each of its bounding rays if and only if its angle exceeds
$\pi$. Near $V$ the region $F^f$ is the wedge bounded by the germs of $E(f,g)$ and $E(f,z_1)$, and $R$
is opposite to the first of those germs, so a point $V+\varepsilon\,d\in R$ with $\varepsilon>0$ small
lies in $F^f$ exactly when the corner at $V$ is reflex.

Now follow $R$ from $V$ to infinity. Far enough out it leaves the bounded region $F^f$; membership in
$F^f$ changes exactly at transverse crossings of $\partial F^f$, and by (G) the ray meets no vertex of
the polygon --- a vertex of $\partial F^f$ on $L_{fg}$ would make four planes concurrent --- and meets
no edge non-transversally except the edge $E(f,g)$ itself, which lies on $L_{fg}$ and on the far side
of $V$. Hence the membership bit at $V+\varepsilon d$ is the parity of the number of edges of
$\partial F^f$ that $R$ crosses.

Which edges are those? Edge $E(f,g)$ is behind $V$. Edge $E(f,z_1)$ has $V$ as an endpoint and its
supporting line meets $L_{fg}$ only at $V$, which is not on the open ray. Edge $E(f,z_0)$ has the
other endpoint $P_{fgz_0}$ of $E(f,g)$ as an endpoint, and its supporting line meets $L_{fg}$ exactly
there, again not on $R$. For every other $v$, the edge $E(f,v)$ lies on $L_{fv}$, which is not parallel to $L_{fg}$ --- that
would make $h_f,h_g,h_v$ share a direction and $\bra{0,f,g,v}$ vanish, against (G) --- and therefore
meets $L_{fg}$ in exactly the point $P_{fgv}$, and that point belongs to $R$ exactly when $B_v=1$ and to the relative interior
of $E(f,v)$ exactly when $C_v=1$ (Lemma~\ref{lem:onb}). At such a point the two supporting lines
are distinct, so the intersection is transverse; by (G) it lies in the interior of a single
boundary edge, not at a polygon vertex. Thus $R$ crosses $E(f,v)$ if and only if
$C_v\wedge B_v=1$, and the parity of the crossings is $\bigoplus_v C_v\wedge B_v$.
\end{proof}

\begin{remark}
Lemma~\ref{lem:parity} is what removes the last free bit from the encoding. By Lemma~\ref{lem:sf} the
reflexness of a corner is a bracket expression \emph{up to} the unknown $s_f$; Lemma~\ref{lem:parity}
computes one particular corner's reflexness from brackets alone, which determines $s_f$ and hence all
eleven corners of that face. Without it, a formula could satisfy every other layer by flipping a
face's orientation bit, and the encoded condition would be strictly weaker.
\end{remark}

\section{The chirotope layers \texorpdfstring{$S$, $H$ and $A$}{S, H and A}}\label{sec:chirotope}

\subsection{Contractions and the cut order}
For a pair $\{i,j\}\subseteq\{0,1,\dots,n\}$ let $\chi_{ij}(x,y)=X_{ijxy}$ on the remaining $n-1$
elements. Each $\chi_{ij}$ is a \emph{tournament}: $\chi_{ij}(y,x)=1\xor\chi_{ij}(x,y)$. There are
$\binom{n+1}{2}=78$ such contractions for $n=12$. The three-term Grassmann--Plücker sign condition for
the rank-$4$ bracket function on a $4$-subset $\{a,b,c,d\}$ of the complement of $\{i,j\}$ is exactly
the rank-$2$ condition for $\chi_{ij}$ on that subset:

\begin{quote}
\textbf{(GP)} the three products $\bra{ijab}\bra{ijcd}$, $-\bra{ijac}\bra{ijbd}$, $\bra{ijad}\bra{ijbc}$
do not all have the same sign.
\end{quote}

This is the content of the Plücker relation
$\bra{ijab}\bra{ijcd}-\bra{ijac}\bra{ijbd}+\bra{ijad}\bra{ijbc}=0$: three non-zero reals summing to
zero cannot share a sign.

\begin{lemma}[Cut order]\label{lem:cutorder}
Let $\chi$ be a tournament on a finite set $X$ with $|X|\ge4$, let $r\in X$, and define for
$x,y\in X\setminus\{r\}$
\[
    d(x,y)\;=\;\chi(x,y)\xor\chi(r,x)\xor\chi(r,y).
\]
Then $\chi$ satisfies \textup{(GP)} on every $4$-subset of $X$ if and only if $d$ is a strict total
order on $X\setminus\{r\}$.
\end{lemma}

\begin{proof}
Reorient every $x\ne r$ by $t_x=1\xor\chi(r,x)$, leaving $r$ unchanged; that is, replace $\chi$ by
$\chi'(u,v)=\chi(u,v)\xor t_u\xor t_v$ with $t_r=0$. Then $\chi'(r,x)=1$ for all $x\ne r$ and
$\chi'(x,y)=d(x,y)$. Reorientation multiplies each of the three products in (GP) by the same sign, so
it preserves (GP).

Suppose (GP) holds. Apply it to a $4$-subset $\{r,x,y,z\}$: the sign bits of the three products are
$1\xor d(y,z)$, $d(x,z)$ and $1\xor d(x,y)$, and they are all equal exactly when
$d(x,y)=d(y,z)=1\xor d(x,z)$, which is precisely one of the two directed $3$-cycles on $\{x,y,z\}$.
Hence $d$ has no directed $3$-cycle. A tournament without directed $3$-cycles is transitive --- if
$x<y$ and $y<z$ but not $x<z$, then $x,y,z$ form one --- so $d$ is a strict total order.

Conversely, suppose $d$ is a strict total order. Choose reals $t_x$ strictly decreasing along it and
put $v_r=(0,1)$ and $v_x=(1,t_x)$ in $\RR^2$. Then $\dett(v_r,v_x)=-1<0$ and
$\dett(v_x,v_y)=t_y-t_x$, so the negativity bits of these $2\times2$ determinants are exactly
$\chi'$. Two-by-two determinants of vectors in the plane satisfy the Plücker identity, so (GP) holds
for $\chi'$ on every $4$-subset, hence for $\chi$.
\end{proof}

\begin{remark}[Why transitivity of the raw tournament would be wrong]\label{rem:counts}
It is tempting to impose transitivity on $\chi_{ij}$ itself. That is false: eleven vectors in a plane
have a \emph{circular} order, and three at $120^{\circ}$ form a directed $3$-cycle. Among the
$2^{\binom{m}{2}}$ tournaments on $m$ elements, exactly $(m-1)!\,2^{\,m-1}$ satisfy (GP) on all
$4$-subsets while only $m!$ are transitive; the ratio is $2^{\,m-1}/m$, which at $m=11$ is
$1024/11\approx93$. Imposing transitivity would delete realizations. Cutting the circle at one element
--- passing from $\chi$ to $d$ --- is what turns the circular order into a linear one, and
Lemma~\ref{lem:cutorder} says nothing is lost or gained by doing so.
\end{remark}

\subsection{The three layers}
Fix a neighbourly map $T$ on $n=12$ faces. The formula has one Boolean variable $X_S$ for each of the
$\binom{13}{4}=715$ brackets, together with the auxiliary variables introduced below.

\begin{description}[leftmargin=2.4em,style=nextline]
\item[Layer $S$ (definitions)] For each of the $66$ pairs $\{i,j\}\subseteq\{1,\dots,12\}$ and each
  pair $\{k,l\}$ of the ten remaining faces, a variable $o^{ij}_{kl}$ with the defining clauses of
  \[
     o^{ij}_{kl}\;=\;X_{0ijk}\xor X_{0ijl}\xor X_{ijkl}.
  \]
  By Theorem~\ref{thm:order} this is the order of $P_{ijk}$ and $P_{ijl}$ along $L_{ij}$; by
  Remark~\ref{rem:cut} it is the cut order of the contraction $\chi_{ij}$ with cut point $0$. Being a
  definition, this layer excludes no assignment of the bracket bits.
\item[Layer $H$ (transitivity on the $66$ affine lines)] For each such pair and each triple $k,l,m$ of
  the remaining faces, the two clauses expressing
  $o^{ij}_{kl}\wedge o^{ij}_{lm}\rightarrow o^{ij}_{km}$ and its converse orientation.
\item[Layer $A$ (transitivity on the $12$ lines at infinity)] For each face $b$ let $r_b$ be the
  smallest face other than $b$. For each pair $\{x,y\}$ of the ten faces other than $b,r_b$ a variable
  \[
      d^{b}_{xy}\;=\;X_{0bxy}\xor X_{0b r_b x}\xor X_{0b r_b y},
  \]
  and for each triple the two transitivity clauses.
\end{description}

Layer $A$ needs the separate treatment because for a pair $\{0,b\}$ the element $0$ is already one of
the two defining labels and cannot serve as the cut point; a different cut point $r_b$ is chosen.
Geometrically, $\bra{0,b,x,y}=-\dett(n_b,n_x,n_y)$ by Lemma~\ref{lem:cofactor}, so the contraction at
$\{0,b\}$ is the chirotope of the eleven other normals modulo $n_b$: eleven vectors in a plane, and
$d^b$ is their cut order.

\begin{theorem}[What $S$, $H$ and $A$ say]\label{thm:chirotope}
An assignment of the $715$ bits $X_S$ extends to a satisfying assignment of $S\wedge H\wedge A$ if and
only if the map $S\mapsto(-1)^{X_S}$, extended by alternation and set to zero on tuples with
repeated labels, is a uniform
chirotope of rank $4$ on the $13$ elements $\{0,1,\dots,12\}$. Moreover every realization in general
position gives such an assignment.
\end{theorem}

\begin{proof}
The auxiliary variables of $S$ and $A$ are defined by their layers, so an assignment of the $X_S$
extends in exactly one way; what $H$ and $A$ then require is that, for each of the $78$ contractions,
the cut order derived from $\chi_{ij}$ --- with cut point $0$ for the $66$ pairs of finite labels and
cut point $r_b$ for the $12$ pairs containing $0$ --- is a strict total order. By
Lemma~\ref{lem:cutorder} this holds for a given contraction if and only if that contraction satisfies
(GP) on all its $4$-subsets. Ranging over the contractions, $S\wedge H\wedge A$ is satisfiable over a
given bracket assignment if and only if all three-term Grassmann--Plücker sign conditions hold.
The alternating extension has every $4$-subset as a basis, so its support is the uniform matroid
$U_{4,13}$. This supplies the matroid-support hypothesis (B1$'$) of
\cite[Theorem 3.6.2, p.~138]{BLSWZ1999}; it must not be omitted for sign maps with zeros.
With $(x_1,x_2,x_3,x_4,y_1,y_2)=(a,b,i,j,c,d)$, condition (B2$''$) of that theorem excludes
the all-negative triple of products in (GP), and interchanging $a,b$ excludes the all-positive
triple. Conversely, (GP) implies this condition; repeated-label instances follow from alternation.
Thus the theorem identifies precisely these conditions with the rank-$4$ chirotope axioms,
proving the first claim.

For the second, let $q$ be a realization in general position. Its brackets are the $4\times4$ minors
of the $13\times4$ matrix of covectors, all non-zero by (G), so they form a uniform rank-$4$ chirotope
and satisfy every Plücker relation. Hence each contraction satisfies (GP), and
Lemma~\ref{lem:cutorder} gives a strict total order for every chosen cut point.
For the finite pairs this is the order of Theorem~\ref{thm:order}, satisfying $S$ and $H$;
for the pairs $\{0,b\}$ it is the cut order $d^b$, satisfying $A$.
This necessary direction uses only determinant identities and Lemma~\ref{lem:cutorder},
not the abstract chirotope characterization of \cite{BLSWZ1999}.
\end{proof}

\begin{remark}
The relational form of the same condition --- one clause of width six per Plücker relation and per
sign pattern --- would need $411\,840$ clauses for $n=12$. The formulation through $78$ cut orders
uses $46\,800$ clauses of width three and four (Appendix~\ref{app:numbers}).
This reduces the clause count; its effect on solving time requires measurement.
The broader use of SAT to study oriented-matroid realizability has precedents in
\cite{Schewe2010,Scheucher2023}. Those references motivate the SAT approach; the
cut-order encoding used here is established by the preceding lemmas.
\end{remark}

\section{The geometric layers \texorpdfstring{$B$ and $I$}{B and I}}\label{sec:geometric}

\subsection{Layer \texorpdfstring{$B$}{B}: face boundaries are simple}
For a face $f$ and two faces $a,b$ that are \emph{not} adjacent in $\link(f)$ --- equivalently, by
Lemma~\ref{lem:adjacency}, with $\{f,a,b\}\notin T$ --- the edges $E(f,a)$ and $E(f,b)$ are non-consecutive
edges of the boundary polygon of $F^f$, and their supporting lines meet at the single point $P_{fab}$.
Layer $B$ consists of the clauses
\[
    \neg\,\OnB(f,a;b)\;\vee\;\neg\,\OnB(f,b;a)
\]
for all such triples, where $\OnB$ is as in Definition~\ref{def:onb}, together with the defining
clauses of the $\OnB$ variables.

\begin{lemma}\label{lem:B}
Every realization in general position satisfies Layer $B$.
\end{lemma}

\begin{proof}
If both bits were $1$ then, by Lemma~\ref{lem:onb}, $P_{fab}$ would be interior to both $E(f,a)$ and
$E(f,b)$, so two non-consecutive closed edges of the boundary polygon of $F^f$ would meet, and the
polygon would not be simple, contradicting (R2).
\end{proof}

\subsection{Layer \texorpdfstring{$I$}{I}: orientation and reflex corners}
Layer $I$ has two parts. For each face $f$ it contains the defining clauses of the corner bits
$\refl(f;a,b)$ of Lemma~\ref{lem:sf} --- one auxiliary variable $s_f$ per face and, per corner, one
turn variable and the one helper variable that splitting its four-term exclusive-or introduces
(Appendix~\ref{app:numbers}) --- together with the \emph{orientation equation} of Lemma~\ref{lem:parity},
\[
   \refl(f;g,z_1)\;=\;\bigoplus_{v\notin\{f,g,z_0,z_1\}} \big(C_v\wedge B_v\big),
\]
where $g$ is the first face of $\link(f)$ and $\{z_0,z_1\}=\Ends(f,g)$.
Either element of $\Ends(f,g)$ may play the role of $z_1$ --- Lemma~\ref{lem:parity} holds for
both, with $B_v$ read relative to the chosen one --- while the corner bit on the left, being attached
to the corner at $P_{f,g,z_1}$, is written with its pair normalized to link-cycle order as in
Lemma~\ref{lem:sf} and Remark~\ref{rem:order-matters}. The two choices give different clause sets, so
one is fixed once and for all, like the emission order, and belongs to the formula's identity; the
implementation takes $z_1$ to be the larger of the two labels. Layer $I$ contains further,
for each of the $44$ vertices $\{f,g,h\}\in T$, the three clauses
\[
   \neg\refl(f;g,h)\vee\neg\refl(g;f,h),\quad
   \neg\refl(f;g,h)\vee\neg\refl(h;f,g),\quad
   \neg\refl(g;f,h)\vee\neg\refl(h;f,g),
\]
each corner bit being written with its pair in link-cycle order as in Remark~\ref{rem:order-matters}.

\begin{lemma}\label{lem:I}
Every realization in general position satisfies Layer $I$, with $s_f$ given by
Lemma~\ref{lem:sf}.
\end{lemma}

\begin{proof}
With $s_f$ the traversal bit of Lemma~\ref{lem:sf}, the variable $\refl(f;a,b)$ takes the value ``the
corner of $F^f$ at $P_{fab}$ is reflex''. The orientation equation then holds by
Lemma~\ref{lem:parity}, and the three clauses at each vertex hold by Lemma~\ref{lem:onereflex}.
\end{proof}

\begin{remark}
The vertex clauses couple the corners of three different faces. They encode the local
intersection restriction of Lemma~\ref{lem:onereflex}; the other geometric clauses concern
one face polygon at a time.
\end{remark}

\section{The gauge}\label{sec:gauge}

The bracket bits of a realization are not unique: two operations change them without changing the
polyhedron.

\begin{lemma}[The reorientation group]\label{lem:gaugegroup}
Let $G=\FF^{\,n+1}$ with coordinates $(\gamma,\epsilon_1,\dots,\epsilon_n)$ act on sign vectors by
\[
   X_S\;\longmapsto\;X_S\;\xor\;\Big(\gamma\;\xor\bigoplus_{i\in S,\ i\ge1}\epsilon_i\Big).
\]
Then $G$ maps the bracket vector of a realization to the bracket vector of a realization of the same
map, and the $\binom{n+1}{4}$ functionals $\varphi_S(\gamma,\epsilon)=\gamma\xor\bigoplus_{i\in S,
i\ge1}\epsilon_i$ span the dual of $G$.
\end{lemma}

\begin{proof}
Replacing $q_i$ by $-q_i$ changes no plane, hence no realization, and flips exactly the brackets
containing $i$: this is $\epsilon_i$. For $\gamma$, let $\tau(x)=Ax+b$ be an orientation-reversing
affine map of $\RR^3$ and $M=\left(\begin{smallmatrix}A&b\\0&1\end{smallmatrix}\right)$, so that a covector $q$
is carried to $qM^{-1}$. The image of a realization is a realization of the same map, since $\tau$ is
a homeomorphism of $\RR^3$ carrying planes to planes. All four rows transform by the same matrix, so
every bracket is multiplied by $\dett(M^{-1})=1/\dett A<0$ and every bit flips. Moreover the last row
of $M^{-1}$ is $(0,0,0,1)$, so $q_0M^{-1}=q_0$: the plane at infinity is fixed, as the action above
assumes.

For the span, suppose a vector $(\gamma,\epsilon)$ is annihilated by every $\varphi_S$, i.e.\
$\varphi_S(\gamma,\epsilon)=0$ for all $S$. Comparing $S=\{i,j,k,l\}$ with $i,j,k,l\ge1$ and
$S'=\{0,j,k,l\}$, which exists because $n\ge4$, gives $\epsilon_i=0$; as $i$ was arbitrary all
$\epsilon_i$ vanish, and then $\gamma=\varphi_S(\gamma,0)=0$. Hence the functionals span the dual.
\end{proof}

\begin{corollary}[Gauge]\label{cor:gauge}
There are brackets $S_k$, for $1\le k\le n+1$, whose functionals $\varphi_{S_k}$
form a basis of the dual of $G$. The action of $G$ on their sign patterns is
simply transitive. Thus every realization can be moved by a unique element of $G$
to one in which all these brackets are positive. Adding a unit clause
$\neg X_{S_k}$ for each $k=1,\dots,n+1$ therefore preserves the property that a realization yields a
satisfying assignment.
\end{corollary}

\begin{proof}
A spanning set contains a basis, which gives the $S_k$; a basis of the dual of a finite-dimensional
$\FF$-vector space makes the action of $G$ on the vector of values simply transitive. Given a
realization, apply the unique group element taking its pattern on $S_1,\dots,S_{n+1}$ to $0$; by
Lemma~\ref{lem:gaugegroup} the result is again a realization of the same map, now satisfying the unit
clauses.
\end{proof}

\begin{remark}
The gauge is a symmetry quotient, not a geometric restriction: it removes $2^{n+1}$-fold redundancy
from the search space. In the implementation the basis is chosen greedily among the brackets that
contain a designated \emph{gauge face} $\varphi$, completed by one bracket that does not contain it.
The completion is necessary, not cosmetic: the functionals of the $\binom{n}{3}$ brackets through a
fixed element $\varphi$ all annihilate $(\gamma,\epsilon)=(1,\epsilon_\varphi{=}1,0,\dots,0)$, so they
span only an $n$-dimensional space and can never form a basis. Twelve of the thirteen pinned brackets
may contain $\varphi$; the thirteenth may not. Which basis is chosen affects the shape of the search
but not the theorem. The gauge is recorded as part of the
formula's identity, because a formula with a different gauge is a different formula.
\end{remark}

\section{The formula and the reduction}\label{sec:formula}

\begin{definition}[The formula]\label{def:formula}
For a neighbourly map $T$ on $n=12$ faces and a gauge face $\varphi$, let $\Phi_T$ --- the gauge face
being fixed once and for all and suppressed from the notation --- be the CNF formula
in the variables
\[
  \underbrace{X_S}_{715}\;,\quad
  \underbrace{o^{ij}_{kl}}_{2\,970}\;,\quad
  \underbrace{d^{b}_{xy}}_{540}\;,\quad
  \underbrace{\OnB(i,j;u)}_{528}\;,\quad
  \underbrace{s_f,\ \text{turn and parity auxiliaries}}_{504}
\]
consisting of the layers $S$, $H$, $A$ (Section~\ref{sec:chirotope}), $B$ and $I$
(Section~\ref{sec:geometric}) emitted in that order, followed by the $13$ gauge clauses of
Corollary~\ref{cor:gauge}. Its size is
\[
\begin{array}{lrr}
\toprule
\text{layer} & \text{clauses} & \text{variables}\\
\midrule
\text{brackets} & 0 & 715\\
S & 23\,760 & 2\,970\\
H & 15\,840 & 0\\
A & 7\,200 & 540\\
B & 2\,640 & 528\\
I & 2\,724 & 504\\
\text{gauge} & 13 & 0\\
\midrule
\Phi_T & 52\,177 & 5\,257\\
\bottomrule
\end{array}
\]
\end{definition}

The emission order is part of the formula's identity: it fixes the numbering of the variables, and a
case or a cube in Section~\ref{sub:cases} names a region of the search space only relative to a fixed
numbering.

\begin{theorem}[Necessity]\label{thm:necessity}
Let $T$ be a neighbourly map on twelve faces and let $q$ be a realization of $T$. Then $\Phi_T$ is
satisfiable.
\end{theorem}

\begin{proof}
By Theorem~\ref{thm:genpos} we may assume $q$ satisfies (G), and by Corollary~\ref{cor:gauge} we may
assume in addition that the $13$ gauge brackets are positive, both replacements producing another
realization of the same $T$. Assign to $X_S$ the negativity bit of $\bra{S}$, to $o^{ij}_{kl}$ and
$d^b_{xy}$ the values forced by their defining clauses, to $\OnB(i,j;u)$ the value forced by
Definition~\ref{def:onb}, to $s_f$ the traversal bit of Lemma~\ref{lem:sf} and to each turn and parity
auxiliary the value forced by its definition. Layer $S$ and the definitional clauses of $A$, $B$ and
$I$ hold by construction; layers $H$ and $A$ hold by Theorem~\ref{thm:chirotope}; layer $B$ by
Lemma~\ref{lem:B}; layer $I$ by Lemma~\ref{lem:I}; the gauge clauses by the choice of representative.
\end{proof}

\begin{corollary}[Conditional non-existence]\label{cor:conditional}
Suppose that for each of the $59$ neighbourly maps $T$ of Theorem~\ref{thm:59} the formula $\Phi_T$ is
unsatisfiable. Then there is no neighbourly polyhedron with twelve faces satisfying
\textup{(H1)}--\textup{(H3)}; that is, the Question of Section~\ref{sec:intro} has a negative answer
within these conventions.
\end{corollary}

\begin{proof}
Such a polyhedron has trivalent vertices (Lemma~\ref{lem:trivalent}) and its induced map is
isomorphic to one of the $59$ representatives $T$ (Theorem~\ref{thm:59}). Relabel its faces by
that isomorphism, so that its induced map is exactly $T$. Lemma~\ref{lem:bridge} then gives a
realization of this representative, and Theorem~\ref{thm:necessity} makes $\Phi_T$ satisfiable,
contrary to hypothesis. Only the necessary direction of Theorem~\ref{thm:chirotope} is used;
this corollary therefore does not depend on the abstract characterization in \cite{BLSWZ1999}.
\end{proof}

\begin{remark}[No realizability conclusion]\label{rem:noconverse}
Satisfiability of some $\Phi_T$ would \emph{not} produce a polyhedron. Two gaps separate a model from
a realization: the layers are necessary conditions and are not claimed to be sufficient, and even a
complete uniform chirotope need not be realizable by real vectors --- non-realizable uniform oriented
matroids exist (see \cite[Sections~2--3]{Scheucher2023} for examples and discussion,
and \cite{BokowskiSturmfels1989} for general realizability methods).
This does not assert that any of our $59$ formulas has a non-realizable model.
A model would be a candidate to be examined, not a geometric counterexample.
\end{remark}

\subsection{Case decomposition and cube trees}\label{sub:cases}
Case decomposition permits parallel execution and bounded, resumable attempts.
It is optional for correctness; the unsplit formula corresponds to depth $d=0$.
We first state the decomposition into the $2^d$ cases of one fixed depth $d$, and then its
generalization to finite trees of cases, which is what the accompanying implementation records.

\begin{proposition}\label{prop:cases}
Let $p_1,\dots,p_d$ be distinct variables of $\Phi_T$ and, for $0\le c<2^{d}$, let $U_c$ be the
conjunction of the $d$ unit clauses assigning $p_m$ the $m$-th bit of $c$. If $\Phi_T\wedge U_c$ is
unsatisfiable for every $c$, then $\Phi_T$ is unsatisfiable.
\end{proposition}

\begin{proof}
A satisfying assignment of $\Phi_T$ satisfies $U_c$ for exactly one $c$, namely the one read off its
values at $p_1,\dots,p_d$.
\end{proof}

A fixed-depth partition can be refined locally: a case that exceeds the available time,
proof size or checker memory can be split again while the other cases are left unchanged.
Cube trees record precisely these successive refinements, so the implementation lets the
tree grow where it is needed. For the rest of this subsection $\Phi$ is any CNF formula in the variables
$1,\dots,V$, and an \emph{assignment} is a map $\alpha\colon\{1,\dots,V\}\to\{0,1\}$; it satisfies
the literal $x$ when $\alpha(x)=1$ and the literal $\neg x$ when $\alpha(x)=0$.

\begin{definition}[Cubes and splits]\label{def:cube}
A \emph{cube} is a set $C$ of literals in which no variable occurs twice; $\varnothing$ is the
\emph{empty cube}. An assignment satisfies $C$ when it satisfies every literal of $C$, and
$\Phi\wedge C$ denotes $\Phi$ together with one unit clause for each literal of $C$. A \emph{split}
of a cube $C$ is a sequence $W=(w_1,\dots,w_k)$, $k\ge1$, of distinct variables none of which
occurs in $C$. Its $2^k$ \emph{children} are the cubes $C\cup\{\ell_1,\dots,\ell_k\}$ with
$\ell_j\in\{w_j,\neg w_j\}$ for each $j$; each of them contains $C$ and has $|C|+k$ literals.
\end{definition}

\begin{lemma}[A split partitions its cube]\label{lem:split}
Let $W$ be a split of a cube $C$. Every assignment that satisfies $C$ satisfies exactly one child
of $C$ under $W$. Hence, if $\Phi\wedge C'$ is unsatisfiable for every child $C'$, then
$\Phi\wedge C$ is unsatisfiable.
\end{lemma}

\begin{proof}
Let $\alpha$ satisfy $C$. Then $\alpha$ satisfies the child $C\cup\{\ell_1,\dots,\ell_k\}$ if and
only if it satisfies each $\ell_j$, and for each $j$ it satisfies exactly one of $w_j$ and
$\neg w_j$. So it satisfies exactly one child: the one with $\ell_j=w_j$ when $\alpha(w_j)=1$ and
$\ell_j=\neg w_j$ when $\alpha(w_j)=0$. If $\alpha$ satisfied $\Phi\wedge C$, it would therefore
satisfy $\Phi\wedge C'$ for that child $C'$.
\end{proof}

\begin{definition}[Cube tree]\label{def:cubetree}
A \emph{cube tree} with root $C_0$ is a rooted tree whose nodes are labelled by cubes, such that
the root is labelled $C_0$ and each node that is not a leaf, say labelled $C$, carries a split $W$
of $C$ and has exactly $2^{|W|}$ children, one for each child of $C$ under $W$ and labelled by it. It
is \emph{refuted} for $\Phi$ when $\Phi\wedge L$ is unsatisfiable for the label $L$ of every leaf.
\end{definition}

A cube tree is finite. Along every path down from the root the labels gain at least one literal
at each step, and no cube has more than $V$ literals, so the height is at most $V-|C_0|$; and each
node has finitely many children.

\begin{theorem}[Coverage]\label{thm:coverage}
Let $\mathcal{T}$ be a cube tree with root $C_0$.
\begin{enumerate}[label=\textup{(\alph*)},leftmargin=2.4em]
\item Every assignment that satisfies $C_0$ satisfies the label of exactly one leaf of
  $\mathcal{T}$.
\item If $\mathcal{T}$ is refuted for $\Phi$, then $\Phi\wedge C_0$ is unsatisfiable. In
  particular, a cube tree with the empty root that is refuted for $\Phi$ shows that $\Phi$ is
  unsatisfiable.
\end{enumerate}
\end{theorem}

\begin{proof}
(a) By induction on the height of $\mathcal{T}$. If the root is a leaf, it is the only leaf and
the claim holds. Otherwise let $W$ be the split at the root, let $C_1,\dots,C_m$, $m=2^{|W|}$, be
the children of $C_0$ under $W$, and let $\mathcal{T}_t$ be the subtree below the child labelled
$C_t$. Each $\mathcal{T}_t$ is a cube tree with root $C_t$ and smaller height, and the leaves of
$\mathcal{T}$ are those of $\mathcal{T}_1,\dots,\mathcal{T}_m$. Let $\alpha$ satisfy $C_0$. By
Lemma~\ref{lem:split}, $\alpha$ satisfies exactly one $C_t$, and by induction it satisfies the
label of exactly one leaf of $\mathcal{T}_t$. For $s\ne t$ every label in $\mathcal{T}_s$ contains
$C_s$, because each label contains the label of its parent; as $\alpha$ does not satisfy $C_s$, it
satisfies no leaf label of $\mathcal{T}_s$.

(b) If an assignment satisfied $\Phi\wedge C_0$, then by (a) it would satisfy the label $L$ of
some leaf, hence $\Phi\wedge L$, which is unsatisfiable by hypothesis. For the empty root,
$\Phi\wedge\varnothing$ is $\Phi$.
\end{proof}

Part (a) says that the leaves of a cube tree divide the region of its root into disjoint regions,
with nothing left over; only part (b) is used below. Proposition~\ref{prop:cases} is the special
case of a tree with at most one split: for $d\ge1$ its cases $U_c$ are the children of the empty
cube under the split $(p_1,\dots,p_d)$, and for $d=0$ the tree consists of its root.

The implementation does not store a tree but everything that was done: possibly several
refutation attempts of one cube, and possibly several splits of one cube recorded by different
machines. The following lemma shows that this is harmless.

\begin{lemma}[Closed cubes]\label{lem:closure}
Let $\mathcal{R}$ be a set of cubes $C$ for which $\Phi\wedge C$ is unsatisfiable, and let
$\mathcal{S}$ be a set of pairs $(C,W)$ of a cube $C$ and a split $W$ of $C$. Define, by induction
on $V-|C|$, a cube $C$ to be \emph{closed} if $C\in\mathcal{R}$, or if there is
$(C,W)\in\mathcal{S}$ such that every child of $C$ under $W$ is closed; the children have more
literals than $C$, so this is already defined for them. Then every closed cube $C$ is the root of
a cube tree, refuted for $\Phi$, whose leaf labels lie in $\mathcal{R}$ and whose splits come from
$\mathcal{S}$. In particular $\Phi\wedge C$ is unsatisfiable for every closed cube $C$.
\end{lemma}

\begin{proof}
By the same induction. If $C\in\mathcal{R}$, the tree with the single node $C$ is refuted.
Otherwise choose $(C,W)\in\mathcal{S}$ whose children are all closed, take for each child a
refuted cube tree with that child as root, which exists by induction, and attach these trees below
a root labelled $C$ carrying the split $W$. The last claim is Theorem~\ref{thm:coverage}(b).
\end{proof}

In the implementation the role of $\mathcal{R}$ is played by the cubes with an admitted record of a
checked refutation, so the lemma applies to the extent that these records are true
(Remark~\ref{rem:attestation}), and that of $\mathcal{S}$ by the recorded splits
(Section~\ref{sec:protocol}). To make its result
reproducible it also selects one tree: at a closed cube it takes a refutation if there is one, and
otherwise the first split in a fixed order whose children are all closed. Every such choice is a
tree as in the proof, and which one is chosen does not affect the conclusion. A cube counts as
closed only through a refutation of that cube or a split of it. The certifier does not
infer closure of $C$ from closure of a proper subset $D\subset C$ of its literals,
although that inference would also be sound because $\Phi\wedge C$ implies $\Phi\wedge D$.

\begin{corollary}[Conditional non-existence from cube trees]\label{cor:treecover}
Suppose that for each of the $59$ neighbourly maps $T$ of Theorem~\ref{thm:59} there is a cube tree
with the empty root that is refuted for $\Phi_T$. Then there is no neighbourly polyhedron with
twelve faces satisfying \textup{(H1)--(H3)}.
\end{corollary}

\begin{proof}
By Theorem~\ref{thm:coverage}(b) each $\Phi_T$ is unsatisfiable, and Corollary~\ref{cor:conditional}
applies.
\end{proof}

\begin{remark}[A closed nonempty root settles only its region]\label{rem:pilotroot}
For a nonempty cube $C_0$, a refuted cube tree with root $C_0$ shows only that no satisfying
assignment of $\Phi$ satisfies $C_0$: a statement about one region of the search space. Pilot
computations are of this kind. A pilot may, for instance, take as its root one case $U_c$ of a
fixed decomposition, written as a cube, and refute a tree below it. That settles one of the $2^d$
cases and says nothing about the others; in particular it does not prove the formula. The
implementation combines closed cubes into a refutation of $\Phi$ only as the leaves of a cube tree
with the empty root --- for example all $2^d$ cases of one decomposition --- and then by
Theorem~\ref{thm:coverage} for that tree. Corollary~\ref{cor:treecover} needs, for every one of
the $59$ formulas, a refuted tree whose root is the empty cube.
\end{remark}

The choice of splitting variables affects cost only, never correctness. The default
campaign first splits the empty cube of each formula on twelve order variables selected
by an adaptation of the proof-prefix heuristic of \cite{BattlemanReevesHeule2025}: bounded
solver runs count occurrences of variables in added clauses. This gives $2^{12}=4\,096$
initial cubes per formula. Subsequent splits use a bounded look-ahead heuristic, with four
variables per split by default, or fewer when the candidate pool or probe budget requires it. The initial depth
can be changed, including to zero; none of these policy choices is a hypothesis of the
reduction. Because Lemma~\ref{lem:split} and
Theorem~\ref{thm:coverage} hold for any choice, the heuristics need no correctness argument of
their own --- but the splits must be recorded, since coverage is a statement about one tree.

\section{The verification protocol}\label{sec:protocol}

A statement of the form ``for each of the $59$ formulas, every leaf of a cube tree with the empty
root is unsatisfiable'' is only as good as the evidence that each leaf was decided, and decided
about the right formula. We state the discipline the accompanying implementation enforces, since
it is part of what would have to be reproduced.

The package \texttt{neighbourly} builds the formulas.
\texttt{cubeproof} runs, records and certifies computations on fixed formulas;
\texttt{faces12-campaign} binds these operations to the $59$ maps and their pinned
formula hashes. A campaign fixes the formulas, requested roots and tool identities.
Workers can process disjoint shards, resume from their journals and transfer open
cubes through recorded handoffs. Ownership determines who works on a region;
only the admitted refutations and splits determine whether that region is closed.
An unfinished or undelivered region remains open.

Three proof modes implement the same acceptance rule. In \texttt{transient} mode
(the default), a proof is written, checked and deleted after its result is committed.
In \texttt{retain} mode, accepted proofs are kept for later checking.
In \texttt{stream} mode, a relay sends proof bytes directly to the checker and hashes
the byte stream; it does not keep a proof file. If a resource limit prevents a check
from completing, the leaf remains open and can be retried or subdivided. These
execution choices do not change $\Phi_T$ or the coverage theorem.

\begin{description}[leftmargin=1.6em,style=nextline]
\item[Two programs, one of them trusted] A CDCL solver (CaDiCaL) is used only as a \emph{proposer}:
  for a leaf cube $L$ it emits an LRAT refutation of $\Phi_T\wedge L$, which \texttt{cake\_lpr}
  \cite{TanHeuleMyreen2023} checks, either from a proof file written once and sealed or streamed
  through a pipe. The correctness of \texttt{cake\_lpr} is a machine-checked theorem carried down to
  machine code by the verified CakeML compiler. A leaf counts as refuted only when the checker's
  standard output is exactly the line \texttt{s VERIFIED UNSAT} followed by a newline,
  the process exits with code zero without a signal, and the CNF it
  read is the one regenerated from the formula and the cube. A leaf that unit propagation alone
  refutes gets a proof generated for it, which passes through the same checker. A solver's own
  claim of unsatisfiability is never accepted.
\item[Identity of the object] A formula is identified by the SHA-256 of its DIMACS text, and the
  bundle that carries it records the map, the layer string and the gauge face from which it was
  generated, so that it can be regenerated and compared byte for byte. A cube is written as its
  literals sorted by variable, so that its identity does not depend on the order in which it was
  reached. The CNF of a leaf is the DIMACS text of the formula, with the clause count in its header
  raised by the number of literals of the cube and one unit clause per literal of the cube appended
  in that order; every record names the hashes of the formula and of this CNF. Evidence about
  different formulas is never merged.
\item[Attestation] Each machine appends what it does to its own journal, a file of JSON records each
  of which carries its own hash and that of its predecessor: refutation attempts with their
  outcomes, splits, and satisfying assignments. A record of a refuted leaf is an
  \emph{attestation}: it states that the checker accepted a proof with a given hash for a CNF with
  a given hash. By default the proof itself is deleted after the check. The hashes detect
  accidental damage and reordering inside a journal; they do not authenticate its writer, and they
  do not detect the replacement of a whole journal or the removal of its last records. The final
  head of each journal that a result relies on must therefore be published independently of the
  file.
\item[Certification] A certifier, which runs neither solver nor checker, reads the journals,
  recomputes the hash of every leaf CNF that a record names from the formula and the cube, admits
  only refutations checked by a checker binary that the campaign names by its hash, and evaluates
  closure as in Lemma~\ref{lem:closure}. A formula counts as proved exactly when its empty cube is
  closed; a closed nonempty root is reported as such (Remark~\ref{rem:pilotroot}). A recorded
  satisfying assignment, which is checked against the clauses, is reported as an inconsistency
  if the empty cube of its formula is closed: by Lemma~\ref{lem:closure} and
  Theorem~\ref{thm:coverage} both cannot be true,
  so one of the attestations or the implementation is at fault. The certificate names the selected tree's refutations, the cubes
  that remain open and one hash over the selected leaves. A second certifier, written independently
  from the specification, must produce the same result from the same journals. Agreement of two
  certifiers checks the evaluation of closure and the selection, not the honesty of the journals.
\item[Independent verification] The default mode does not retain proofs. Their total size
  depends on the eventual tree and solver runs; no storage estimate for a complete
  computation is claimed. An attested leaf is verified again in one of two ways. Where a proof was retained,
  the checker checks it again against the regenerated CNF, without the solver. Otherwise a second
  party recomputes the leaf: it regenerates the CNF and compares its hash with the record, runs the
  recorded solver again (or, for a leaf refuted by unit propagation, generates its proof again),
  has the checker verify the new proof, and compares the new proof's hash with the recorded one. Strict reproduction requires identical proof bytes under the recorded binary, options
  and execution profile; a difference fails that reproduction check. Determinism is tested
  on small workloads, not assumed as a mathematical theorem. A different checked proof
  still refutes the same leaf, but does not reproduce the recorded proof bytes. Either way
  the leaf is established anew, independently of the journal. Verifying a random sample of leaves
  tests the journals statistically and verifies only the leaves in the sample; verifying every
  selected leaf of every formula repeats the computation and verifies it completely. Neither a
  journal nor a proof is accepted without the corresponding check.
\item[Exact controls and a reference encoder] The identities of Section~\ref{sec:dictionary} are
  checked against exact rational arithmetic on explicit plane configurations: random ones in general
  position for Theorem~\ref{thm:order}, Lemma~\ref{lem:onb} and Lemma~\ref{lem:turn}, and, for Layers
  $B$ and $I$, which hold only for realizations, the tetrahedron and an exact rational realization of
  the Szilassi polyhedron \cite{GrunbaumSzilassi2009} perturbed into general position, for which the
  predicates have separately computed truth values. Two reference checks serve different purposes.
  The package reference encoder shares the production encoder's conventions and compares clauses,
  detecting transcription errors. The additional reference in \texttt{tests/reference} is written
  from this article with its own bracket numbering and sign calculations. It reads the emitted
  clauses and the names of their variables, recognizes the auxiliary definitions by their truth
  tables, and compares the resulting Boolean polynomials over $\FF$ with the predicates defined
  here. It also checks the substantive clauses, the gauge rank and reorientation invariance.
  This symbolic comparison covers all assignments of the bracket bits for each of the $59$
  fixed maps; separate rational tests check geometric interpretations on sampled arrangements
  and on the two known realizations, including relabellings and changes of gauge face.
  The reference reads the encoder's choices of link traversal and pinned corner and checks
  their admissibility; the pinned formula hashes fix the exact choices. Both references are
  maintained by this project. Neither constitutes external peer review or proves the geometric
  lemmas from agreement between programs.
\end{description}

The assurance argument has three distinct levels. The geometric reduction and the
completeness of the case list are mathematical obligations. Correct encoding, formula
identity and complete coverage are implementation obligations, supported here by exact
controls and separate computations; coverage itself is Theorem~\ref{thm:coverage}, and the
certifiers only evaluate its hypothesis on recorded evidence. Finally, the checker establishes
unsatisfiability of the concrete CNF it reads, subject to the execution assumptions of its
verification theorem; the small C layer through which the verified machine code reads its input
files, and, when a proof is streamed, the pipe through which it arrives, lie outside that theorem;
they are identified by the hashes recorded with each result. Formal verification of the checker does not discharge the
first two levels. The solver need not be trusted to assert unsatisfiability: its proposed
refutation is checked independently. The resulting theorem would nevertheless depend on the whole
chain from geometry to the checked formulas.

\begin{remark}[Attestation and verification are different guarantees]\label{rem:attestation}
The splits of a cube tree need no trust: whether a recorded split is a split in the sense of
Definition~\ref{def:cube} is decided by reading it, and its children are determined by it. Only the
leaves carry computational content. A certificate that a formula is proved establishes the
hypothesis of Theorem~\ref{thm:coverage}(b) for the selected tree at the level of attestation: each
leaf is covered by a record of an accepted check. It does not show that the writer of an admitted
record actually ran the checker on that CNF; a journal is not a proof. That a leaf is
unsatisfiable is established independently of the journals only by checking a retained proof
again or by recomputing the leaf. An independent verification of the whole computation therefore
verifies every selected leaf of every one of the $59$ formulas; a sample verifies only the leaves
it selects.
\end{remark}

\section{Status, scope and what remains}\label{sec:status}

\subsection{What has been done}
The geometric and encoding lemmas, the reduction and the coverage statements are
proved here. The case-list theorem uses the published classification, whose concrete
representatives are also recomputed and compared with the published data. The enumeration of
Theorem~\ref{thm:59} has been carried out and agrees in number with
\cite{AltshulerBokowskiSchuchert1996} and class by class with \cite{LutzManifoldPage}. The encoder produces the $59$ formulas of
Definition~\ref{def:formula}, and the finite identities are tested by machine:
Theorem~\ref{thm:order}, Lemma~\ref{lem:onb} and Lemma~\ref{lem:turn} against exact arithmetic on
random rational configurations, Lemma~\ref{lem:cutorder} exhaustively on all tournaments with at most
seven elements, which also reproduces the counts of Appendix~\ref{app:numbers},
Lemma~\ref{lem:onereflex} through the supplementary check of Remark~\ref{rem:octants}, and
the gauge of Corollary~\ref{cor:gauge} by running all $2^{13}$ group elements and observing that they
give the thirteen pinned brackets all $8\,192$ possible sign patterns within each orbit. The gauge
selects one representative from each orbit; it does not identify the $8\,192$ patterns with
different orbits. The tests of Lemmas~\ref{lem:onb} and~\ref{lem:turn} evaluate the encoder's own
literals; the others check the identities independently of it. None of them replaces a proof given
here.

\subsection{What has not been done}
A complete set of checked refutations has not been obtained: for none of the $59$ formulas has a
refuted cube tree with the empty root been recorded, let alone verified independently.
Corollaries~\ref{cor:conditional} and~\ref{cor:treecover} therefore remain conditional: this paper
does not settle the Question. Bounded pilot computations on restricted roots are not partial
proofs of it beyond the regions they close (Remark~\ref{rem:pilotroot}). This version is a preprint
of the reduction and of the protocol, released before the computation; the computation, and its
independent verification, will be reported in a revised version. Clause counts and bounded tests do
not provide a reliable estimate of the cost of all $59$ instances. Cube trees support parallel and
resumable attempts, but timeouts, resource limits and unrefuted leaves must remain explicitly open
in any campaign report.

\subsection{Scope}
The reduction is stated for polyhedra satisfying (H1)--(H3). Dropping (H1) to allow two faces to meet
in more than one edge or in an isolated additional point, dropping (H2) to allow faces with holes, or
allowing self-intersections, changes the object: vertices of degree exceeding three become possible,
the case list of Theorem~\ref{thm:59} no longer applies, and the $59$ formulas say nothing about those
variants. Statements made here for general $n$, such as Proposition~\ref{prop:euler} and
Lemma~\ref{lem:framing}, are proved; the computations, and the tests of the implementation on maps
with other numbers of faces, are finite checks of particular instances and establish nothing about
polyhedra with other numbers of faces. We also note again the difference from the dual problem: Schewe's theorem
\cite{Schewe2010} excludes polyhedral realizations of the $12$-vertex genus-$6$ triangulations, which
are combinatorially dual to our maps; realizability is not preserved by that duality, so it does not transfer
a non-realizability result from one problem to the other.

\appendix

\section{Numerical data}\label{app:numbers}

\subsection{The case list}
The enumeration of Theorem~\ref{thm:59} returns, for $n=12$ with the link of one face fixed,
$12\,144$ labelled orientable maps in $59$ isomorphism classes. Their automorphism orders are
\[
\begin{array}{lrrrrrr}
\toprule
|\Aut| & 1 & 2 & 3 & 4 & 6 & 12\\
\text{classes} & 39 & 9 & 1 & 7 & 2 & 1\\
\bottomrule
\end{array}
\qquad
\sum_{C}\frac{264}{|\Aut(C)|}=12\,144 .
\]
On the smaller sizes the same program returns one class for $n=4$ ($|\Aut|=24$), one for
$n=7$ ($|\Aut|=42$, the Császár torus) and none for $n=10$.

The automorphism orders are those stated in \cite{EllinghamStephens2005}, Theorem~4. The $59$
classes are those of the list of \cite{LutzManifoldPage} for genus $6$ and twelve vertices. The full correspondence, including all twelve-label bijections, is stored in
\texttt{data/fixtures/lutz\_correspondence.json}. The test applies them and compares
the normalized file hash with the published file; the source URL, retrieval date
and both raw and normalized hashes are recorded in \texttt{docs/SOURCES.md}.
Keeping the correspondence with the data avoids a second, manually maintained
numbering table in the article.

\subsection{Tournaments}
The counts behind Remark~\ref{rem:counts}, obtained by exhaustive enumeration of all tournaments
(the accompanying test suite recomputes them):
\[
\begin{array}{lrrrr}
\toprule
m & 4 & 5 & 6 & 7\\
\midrule
\text{tournaments satisfying (GP)} & 48 & 384 & 3\,840 & 46\,080\\
\text{transitive tournaments} & 24 & 120 & 720 & 5\,040\\
\bottomrule
\end{array}
\]
matching $(m-1)!\,2^{\,m-1}$ and $m!$ respectively.

\subsection{The formula}
Sizes are as in Definition~\ref{def:formula}; the distribution of clause widths, counted from the
emitted formula of the first map, is
\[
\begin{array}{lrrrrr}
\toprule
\text{width} & 1 & 2 & 3 & 4 & \text{total}\\
\midrule
S      &   &     & & 23\,760 & 23\,760\\
H      &   &     & 15\,840 & & 15\,840\\
A      &   &     & 2\,880 & 4\,320 & 7\,200\\
B      &   & 528 & 2\,112 & & 2\,640\\
I      &   & 324 & 1\,056 & 1\,344 & 2\,724\\
\text{gauge} & 13 & & & & 13\\
\midrule
\Phi_T & 13 & 852 & 21\,888 & 29\,424 & 52\,177\\
\bottomrule
\end{array}
\]
The pattern behind the table: a three-term exclusive-or definition contributes eight clauses of width
four, a two-term one contributes four clauses of width three, a conjunction contributes one clause of
width three and two of width two, transitivity contributes clauses of width three, and the boundary
clauses of $B$ and the vertex clauses of $I$ have width two. Exclusive-ors of more than three terms,
which occur only in the orientation equation and in the four-term exclusive-or that defines each
corner bit, are split into three-term ones with fresh variables; any assignment of the inputs extends uniquely to the fresh variables, so
splitting changes neither the models nor the necessity argument. For the first map of the shipped list
the DIMACS text has $5\,257$ variables and $52\,177$ clauses.

\section{The accompanying software}\label{app:software}

The software is published in the repository
\begin{center}
\url{https://github.com/pirozhkov-andrei/polyhedraproof}
\end{center}
This revision of the article was released with version 1.0.0rc1 of the software, under the tag
\texttt{v1.0.0rc1}; the source archive of that release holds the article, the code, the data and the
tests referred to here. The implementation is written in Python, needs nothing beyond the standard
library at run time,
and runs CaDiCaL and \texttt{cake\_lpr} as external programs. Its code
refers to the statements of this paper by their \LaTeX{}
labels --- \texttt{lem:parity} for Lemma~\ref{lem:parity}, for instance --- which do not change when
the numbering does. The correspondence itself, statement by statement, with the code that implements
each one and the tests that check it, is the file \texttt{docs/THEORY\_MAP.md} of the repository. It
is generated from this paper's source and from a register of claims kept with the code, and a test
fails when it is out of date or names a statement, a function or a test that does not exist; so it
is kept there rather than copied here, where it could only go stale.

\section*{Use of generative artificial intelligence}
Generative AI tools assisted with reviewing mathematical arguments, developing and checking the
accompanying code, and drafting and revising the English exposition. These activities were part of
the development process; AI-generated suggestions and reviews are not evidence of mathematical
correctness, independent peer review, or formal verification. The author takes responsibility for
the mathematical claims, the software, and the final text. The verification protocol described in
Section~\ref{sec:protocol} uses the published source code and the LRAT checker, not AI judgements.

\begin{thebibliography}{99}

\bibitem{AltshulerBokowskiSchuchert1996}
A.~Altshuler, J.~Bokowski and P.~Schuchert,
\emph{Neighborly 2-manifolds with 12 vertices},
J. Combin. Theory Ser. A \textbf{75} (1996), no.~1, 148--162.
\href{https://doi.org/10.1006/jcta.1996.0069}{doi:10.1006/jcta.1996.0069}.

\bibitem{ArsenevaEtAl2024}
E.~Arseneva, L.~Kleist, B.~Klemz, M.~Löffler, A.~Schulz, B.~Vogtenhuber and A.~Wolff,
\emph{Adjacency graphs of polyhedral surfaces},
Discrete Comput. Geom. \textbf{71} (2024), no.~4, 1429--1455.
\href{https://doi.org/10.1007/s00454-023-00537-6}{doi:10.1007/s00454-023-00537-6};
preliminary version in \emph{Proc. 37th Int. Symp. on Computational Geometry (SoCG 2021)},
LIPIcs \textbf{189}, 11:1--11:17.

\bibitem{BattlemanReevesHeule2025}
Z.~Battleman, J.~E.~Reeves and M.~J.~H.~Heule,
\emph{Problem partitioning via proof prefixes},
in \emph{Proc. 28th Int. Conf. on Theory and Applications of Satisfiability Testing (SAT 2025)},
LIPIcs \textbf{341}, 3:1--3:18.
\href{https://doi.org/10.4230/LIPIcs.SAT.2025.3}{doi:10.4230/LIPIcs.SAT.2025.3}.

\bibitem{BiereEtAl2024}
A.~Biere, T.~Faller, K.~Fazekas, M.~Fleury, N.~Froleyks and F.~Pollitt,
\emph{CaDiCaL 2.0},
in \emph{Computer Aided Verification (CAV 2024)}, LNCS \textbf{14681}, Springer, 2024, 133--152.
\href{https://doi.org/10.1007/978-3-031-65627-9_7}{doi:10.1007/978-3-031-65627-9\_7}.

\bibitem{BLSWZ1999}
A.~Björner, M.~Las~Vergnas, B.~Sturmfels, N.~White and G.~M.~Ziegler,
\emph{Oriented Matroids}, 2nd ed.,
Encyclopedia of Mathematics and its Applications \textbf{46}, Cambridge University Press, 1999.
\href{https://doi.org/10.1017/CBO9780511586507}{doi:10.1017/CBO9780511586507}.

\bibitem{BokowskiKevin2025}
J.~Bokowski and Kevin~H.,
\emph{Polyhedral embeddings of triangular regular maps of genus $g$, $2\le g\le14$,
and neighborly spatial polyhedra},
Symmetry \textbf{17} (2025), no.~4, article~622.
\href{https://doi.org/10.3390/sym17040622}{doi:10.3390/sym17040622}.

\bibitem{BokowskiSturmfels1989}
J.~Bokowski and B.~Sturmfels,
\emph{Computational Synthetic Geometry},
Lecture Notes in Mathematics \textbf{1355}, Springer, 1989.
\href{https://doi.org/10.1007/BFb0089253}{doi:10.1007/BFb0089253}.

\bibitem{EllinghamStephens2005}
M.~N.~Ellingham and C.~Stephens,
\emph{Triangular embeddings of complete graphs (neighborly maps) with 12 and 13 vertices},
J. Combin. Des. \textbf{13} (2005), no.~5, 336--344.
\href{https://doi.org/10.1002/jcd.20048}{doi:10.1002/jcd.20048}.

\bibitem{GrunbaumSzilassi2009}
B.~Grünbaum and L.~Szilassi,
\emph{Geometric realizations of special toroidal complexes},
Contrib. Discrete Math. \textbf{4} (2009), no.~1, 21--39.
\href{https://doi.org/10.55016/ojs/cdm.v4i1.61986}{doi:10.55016/ojs/cdm.v4i1.61986}.

\bibitem{HackerPoet}
HackerPoet (CodeParade), \emph{NeighborlyPolyhedra}, software repository,\newline
\url{https://github.com/HackerPoet/NeighborlyPolyhedra}.
Accessed 15 October 2026. Related published account: \cite{BokowskiKevin2025}.

\bibitem{Hatcher2002}
A.~Hatcher, \emph{Algebraic Topology}, Cambridge University Press, 2002.
Author's online edition: \url{https://pi.math.cornell.edu/~hatcher/AT/ATpage.html}.

\bibitem{JungermanRingel1980}
M.~Jungerman and G.~Ringel,
\emph{Minimal triangulations on orientable surfaces},
Acta Math. \textbf{145} (1980), no.~1, 121--154.
\href{https://doi.org/10.1007/BF02414187}{doi:10.1007/BF02414187}.

\bibitem{Lazarus2020}
F.~Lazarus, \emph{Linear Embeddings}, lecture notes, 2020, Section~2.
\url{https://pagesperso.g-scop.grenoble-inp.fr/~lazarusf/Enseignement/effectiveGeom3.pdf}.

\bibitem{LutzManifoldPage}
F.~H.~Lutz, \emph{The Manifold Page}, section \emph{Surfaces}.
\url{https://www3.math.tu-berlin.de/IfM/Nachrufe/Frank_Lutz/stellar/surfaces.html}.
Retrieved 13 October 2026.

\bibitem{Meisters1975}
G.~H.~Meisters, \emph{Polygons have ears},
Amer. Math. Monthly \textbf{82} (1975), no.~6, 648--651.
\href{https://doi.org/10.1080/00029890.1975.11993898}{doi:10.1080/\allowbreak 00029890.1975.11993898}.

\bibitem{Ringel1974}
G.~Ringel, \emph{Map Color Theorem},
Grundlehren der mathematischen Wissenschaften \textbf{209}, Springer, 1974.
\href{https://doi.org/10.1007/978-3-642-65759-7}{doi:10.1007/978-3-642-65759-7}.

\bibitem{RingelYoungs1968}
G.~Ringel and J.~W.~T.~Youngs,
\emph{Solution of the Heawood map-coloring problem},
Proc. Nat. Acad. Sci. U.S.A. \textbf{60} (1968), no.~2, 438--445.
\href{https://doi.org/10.1073/pnas.60.2.438}{doi:10.1073/pnas.60.2.438}.

\bibitem{Scheucher2023}
M.~Scheucher,
\emph{A SAT attack on Erd\H{o}s--Szekeres numbers in $\RR^d$ and the empty hexagon theorem},
Comput. Geom. Topol. \textbf{2} (2023), no.~1, 2:1--2:13.
\href{https://doi.org/10.57717/cgt.v2i1.12}{doi:10.57717/cgt.v2i1.12}.

\bibitem{Schewe2010}
L.~Schewe,
\emph{Nonrealizable minimal vertex triangulations of surfaces: showing nonrealizability using oriented
matroids and satisfiability solvers},
Discrete Comput. Geom. \textbf{43} (2010), no.~2, 289--302.
\href{https://doi.org/10.1007/s00454-009-9222-y}{doi:10.1007/s00454-009-9222-y};
\href{https://arxiv.org/abs/0801.2582}{arXiv:0801.2582}.

\bibitem{SulankeLutz2009}
T.~Sulanke and F.~H.~Lutz,
\emph{Isomorphism-free lexicographic enumeration of triangulated surfaces and 3-manifolds},
European J. Combin. \textbf{30} (2009), 1965--1979.
\href{https://doi.org/10.1016/j.ejc.2008.12.016}{doi:10.1016/j.ejc.2008.12.016};
\href{https://arxiv.org/abs/math/0610022}{arXiv:math/0610022}.

\bibitem{Szilassi1986}
L.~Szilassi, \emph{Regular toroids}, Structural Topology \textbf{13} (1986), 69--80.

\bibitem{TanHeuleMyreen2023}
Y.~K.~Tan, M.~J.~H.~Heule and M.~O.~Myreen,
\emph{Verified propagation redundancy and compositional UNSAT checking in CakeML},
Int. J. Softw. Tools Technol. Transf. \textbf{25} (2023), no.~2, 167--184.
\href{https://doi.org/10.1007/s10009-022-00690-y}{doi:10.1007/s10009-022-00690-y};
conference version: \emph{cake\_lpr: verified propagation redundancy checking in CakeML},
TACAS 2021, LNCS \textbf{12652}, 223--241.

\end{thebibliography}

\end{document}
